% ============================================================
% preambulo-hyperref.tex
% ------------------------------------------------------------
% CUARTA Y ÚLTIMA CAPA DEL PREÁMBULO. VA DESPUÉS DE LA ESTÉTICA.
%
%   % ============================================================
% preambulo-proyecto.tex
% ------------------------------------------------------------
% ARCHIVO DERIVADO — GENERADO POR gbpublisher
% NO EDITAR A MANO: SE REESCRIBE EN CADA GENERACIÓN DEL
% main.tex. LOS CAMBIOS SE HACEN EN LA BASE.
%
% SÍ SE VERSIONA: EL COMMIT ES LA GARANTÍA DE TRAZA DE QUÉ SE
% PUBLICÓ. POR ESO NO LLEVA MARCA DE TIEMPO.
% ------------------------------------------------------------
% ESTA CAPA SOLO DEFINE VALORES. NO CARGA NINGÚN PAQUETE.
% VA PRIMERA PORQUE geometry Y fontspec NECESITAN LOS VALORES
% EN EL MOMENTO DEL \usepackage.
% ============================================================

% --- PRESTACIONES DE SALIDA (derivadas de estado_libro) ---
\newif\ifgbshowframe    \gbshowframefalse
\newif\ifgbhomeoarchy   \gbhomeoarchytrue
\newif\ifgbenlaces      \gbenlacestrue
\newif\ifgbcruces       \gbcrucesfalse

% --- FORMATO DE PÁGINA (formatos_pdf id=32) ---
\def\gbanchopagina{15.50cm}
\def\gbaltopagina{23.00cm}
\def\gbmargeninterior{2.00cm}
\def\gbmargenexterior{2.00cm}
\def\gbmargensuperior{2.00cm}
\def\gbmargeninferior{2.00cm}
\def\gbnombreformato{Estudio 2A / 04}
\def\gbanchomontaje{18.00truecm}
\def\gbaltomontaje{25.50truecm}

% --- IDIOMA ---
\def\gbidiomas{main=spanish}
\def\gbidiomaprincipal{spanish}
\newif\ifgbcjk   \gbcjkfalse

% --- BIBLIOGRAFÍA ---
\def\estandar{apa}
\newif\ifgbbibporcapitulo   \gbbibporcapitulotrue
\def\gbarchivobib{./referencias/ref-l-DUEK.bib}

% --- METADATOS DEL PDF ---
\def\gbtitulo{Un tiempo diferente}
\def\gbsubtitulo{Vida cotidiana y pandemia}
\def\gbtituloabreviado{}
\def\gbautores{Moguillansky, Marina (comp.); Duek, Carolina (comp.)}
\def\gbeditorial{Ediciones Imago Mundi}
\def\gbciudad{Buenos Aires}
\def\gbanio{2026}
\def\gbisbn{978-950-793-487-2}
\def\gbisbnelectronico{}
\def\gbdoi{}
\def\gblicencia{}
\def\gburllicencia{}
\def\gbpalabrasclave{}

% --- VERSIÓN DE gbpublisher ---
\def\gbversion{0.0.332}
   valores de la base
%   % ============================================================
% preambulo-contrato.tex
% ------------------------------------------------------------
% CONTRATO VERSIÓN 1
% ------------------------------------------------------------
% DEFINE TODO LO QUE docbook-to-latex.xsl PUEDE EMITIR EN LOS
% FRAGMENTOS DE CAPÍTULO. ES LA CAPA QUE UNA EDITORIAL
% CONSERVA CUANDO REEMPLAZA LA ESTÉTICA: SI FALTA UNA DE
% ESTAS DEFINICIONES, EL LIBRO NO COMPILA; SI ESTÁ DEFINIDA
% DISTINTO, COMPILA Y SALE MAL.
%
% CUANDO SE AGREGUE VOCABULARIO NUEVO, SUBIR LA VERSIÓN DEL
% CONTRATO. EL main.tex DECLARA LA QUE NECESITA Y ESTA CAPA
% LA VERIFICA: ASÍ UN PREÁMBULO VIEJO FRENA EN LA PRIMERA
% LÍNEA CON UN MENSAJE ÚTIL, EN VEZ DE DAR UN
% "Undefined control sequence" EN LA PÁGINA 140.
% ------------------------------------------------------------
% DEPENDE DE preambulo-proyecto.tex, QUE VA ANTES.
% ============================================================

\def\gbContratoVersion{1}

% ============================================================
% \gbRequiereContrato{N}
% PROPÓSITO : EL main.tex DECLARA QUÉ VERSIÓN NECESITA.
% PARÁMETRO : N — versión requerida
% ============================================================
\newcommand{\gbRequiereContrato}[1]{%
  \ifnum#1>\gbContratoVersion\relax
    \PackageError{gbpublisher}{%
      Este preambulo implementa el contrato \gbContratoVersion\space
      y el libro requiere el #1}{%
      Actualizar preambulo-contrato.tex.}%
  \fi}

% ============================================================
% HOJAS DE CORTESÍA
% ------------------------------------------------------------
% Dos páginas en blanco al inicio —una hoja— que todo libro
% lleva. Entran en la paginación: son i y ii, y el sumario
% empieza en iii.
%
% \hbox{} ES IMPRESCINDIBLE: UNA PÁGINA SIN NINGÚN MATERIAL NO
% SE EMITE, ASÍ QUE UN \newpage SOBRE UNA PÁGINA VACÍA NO
% PRODUCE NADA Y EL CONTADOR NO AVANZA. LA CAJA VACÍA DA EL
% MÍNIMO DE CONTENIDO QUE OBLIGA A COMPONERLA.
%
% \thispagestyle{empty} LES QUITA EL FOLIO: SE CUENTAN PERO NO
% SE IMPRIME EL NÚMERO.
% ============================================================
\newcommand{\gbHojasCortesia}{%
	\thispagestyle{empty}\hbox{}\newpage
	\thispagestyle{empty}\hbox{}\newpage}

% ============================================================
% ENCABEZADO DE PIEZA GENERADA
% ------------------------------------------------------------
% Para las piezas sin cuerpo que llevan título propio. El título
% sale de capitulos.titulo_capitulo, que es donde el editor lo
% escribe, y no del preámbulo.
%
% \markboth actualiza el folio corrido, que \chapter* no toca.
% ============================================================
\newcommand{\gbEncabezadoPieza}[1]{%
	\chapter*{#1}%
	\addcontentsline{toc}{chapter}{#1}%
	\markboth{\MakeUppercase{#1}}{\MakeUppercase{#1}}}

% ============================================================
% ÍNDICE CON SU ENCABEZADO
% ------------------------------------------------------------
% \printindex HACE UN \clearpage INTERNO AUNQUE EL ÍNDICE SE
% DECLARE CON title={}. SIN NEUTRALIZARLO, EL ENCABEZADO QUEDA
% SOLO EN UNA PÁGINA Y LA LISTA EMPIEZA EN LA SIGUIENTE.
%
% EL \let VA DENTRO DE UN GRUPO: FUERA DE ÉL \clearpage TIENE QUE
% SEGUIR FUNCIONANDO, Y EL COLOFÓN QUE VIENE DESPUÉS LO NECESITA.
% ============================================================
\newcommand{\gbIndice}[2]{%
	\gbEncabezadoPieza{#1}%
	\begingroup
	\let\clearpage\relax
	\raggedright
	\small
	\printindex[#2]%
	\endgroup}

% ============================================================
% 1. PAQUETES QUE EL CONTRATO NECESITA
% ------------------------------------------------------------
% SOLO LO QUE HACE FALTA PARA QUE UN FRAGMENTO COMPILE. LA
% ESTÉTICA CARGA LO SUYO APARTE. TODOS ESTÁN EN TeX Live Y
% NINGUNO PIDE --shell-escape.
% ============================================================
\usepackage{etoolbox}
\usepackage{graphicx}
\usepackage{booktabs}
\usepackage{array}
\usepackage{listings}
\usepackage{ragged2e}
\usepackage{xcolor}
\usepackage[most]{tcolorbox}

% ------------------------------------------------------------
% IDIOMA
% ------------------------------------------------------------
% babel VA EN EL CONTRATO Y NO EN LA ESTÉTICA, POR DOS RAZONES.
% LA PRIMERA ES DE ORDEN: biblatex PIDE CARGARSE DESPUÉS DE
% babel Y csquotes, Y LA ESTÉTICA VA DESPUÉS DE ESTA CAPA.
% LA SEGUNDA ES DE FONDO: EL IDIOMA DECIDE LAS COMILLAS QUE
% PRODUCE \enquote, LA PARTICIÓN DE PALABRAS Y LAS CADENAS DE
% LA BIBLIOGRAFÍA. SON GARANTÍAS DEL CONTRATO, NO DECISIONES
% DE DISEÑO: UNA EDITORIAL QUE REEMPLAZA LA ESTÉTICA NO DEBE
% PODER CAMBIAR EL IDIOMA DEL LIBRO.
\usepackage[\gbidiomas]{babel}
\frenchspacing

% COMILLAS SEGÚN EL IDIOMA ACTIVO: EN ESPAÑOL DA « ».
% EL XSLT EMITE \enquote Y NUNCA COMILLAS LITERALES.
\usepackage[autostyle=true]{csquotes}


% ============================================================
% 3. BIBLIOGRAFÍA
% ------------------------------------------------------------
% EL ARCHIVO biblatex-<estandar>-config.tex ES RESPONSABLE DE
% CARGAR biblatex CON EL ESTILO QUE CORRESPONDA. SON NUESTROS
% ARCHIVOS, VIAJAN EN EL PAQUETE Y SE COPIAN AL PROYECTO: NO
% SE DEPENDE DE NINGÚN ESTILO REGISTRADO EN TeX Live.
%
% LA RAMA bibtex (vancouver) NO ESTÁ IMPLEMENTADA. EL XSLT
% CORTA ANTES DE GENERAR SI EL ESTILO LA REQUIERE.
% ============================================================
\input{cita-\estandar-config.tex}

% ============================================================
% Índice de autores del aparato bibliográfico
% ------------------------------------------------------------
% Lo llena biblatex: con indexing=true en el estilo, cada nombre
% citado y cada nombre de la bibliografía se indexan solos. El
% editor no escribe ningún \index.
%
% esindex ORDENA LOS APELLIDOS COMPUESTOS CASTELLANOS COMO
% CORRESPONDE: «de la Torre» SE ALFABETIZA POR Torre PERO SE
% IMPRIME CON LA PARTÍCULA. SIN ÉL, EL ORDEN POR OMISIÓN LOS
% MANDA A LA D.
%
% VA EN EL CONTRATO PORQUE ES LO QUE HACE QUE LA PIEZA
% indice_autores_bibliografia DIGA LA VERDAD. EL \makeindex ES
% DISEÑO Y VA EN LA ESTÉTICA.
% ============================================================
\usepackage{esindex}
\DeclareIndexNameFormat{default}{%
	\usebibmacro{index:name}{\esindex[names]}
	{\namepartfamily}
	{\namepartgiven}
	{\namepartprefix}
	{\namepartsuffix}}

% ============================================================
% 4. PÁGINAS Y SALTOS
% ============================================================

% ------------------------------------------------------------
% \gbclearrecto — SALTA A LA PRÓXIMA PÁGINA IMPAR
% \gbclearverso — SALTA A LA PRÓXIMA PÁGINA PAR
% ------------------------------------------------------------
% LA PÁGINA DE RELLENO LLEVA \hbox{} PORQUE UNA PÁGINA SIN
% CONTENIDO NO SE EMITE. NO SE USA UN CARÁCTER BLANCO: ESO
% DEJA TEXTO EXTRAÍBLE EN EL PDF, QUE APARECE AL COPIAR, EN
% pdftotext Y EN UN LECTOR DE PANTALLA.
% EL CONTADOR SIGUE CORRIENDO: SOLO SE OCULTA EL FOLIO.
% ------------------------------------------------------------
\makeatletter
\newcommand{\gbclearrecto}{%
  \clearpage
  \ifodd\c@page\else
    \thispagestyle{empty}\hbox{}\newpage
  \fi}
\newcommand{\gbclearverso}{%
  \clearpage
  \ifodd\c@page
    \thispagestyle{empty}\hbox{}\newpage
  \fi}
% LA BLANCA QUE GENERA \cleardoublepage SALE, POR OMISIÓN, CON
% EL FOLIO Y EL ENCABEZADO CORRIENTES. ESTE PARCHE LA DEJA
% VACÍA, QUE ES LO QUE SE ESPERA DE UNA PÁGINA EN BLANCO.
\renewcommand{\cleardoublepage}{\gbclearrecto}
\makeatother

% ------------------------------------------------------------
% gbpaginaespecial / gbpaginaimagen
% ------------------------------------------------------------
% PIEZAS SUELTAS DE PRELIMINARES: UNA DEDICATORIA, UN
% RECORDATORIO, UNA FOTO A PÁGINA ENTERA. VAN EN RECTO, SIN
% FOLIO NI ENCABEZADO, Y SU VUELTA QUEDA EN BLANCO. NO LLEVAN
% \chapter Y NO ENTRAN AL ÍNDICE.
% LA DE IMAGEN ROMPE LA CAJA DE TEXTO; LA OTRA LA RESPETA.
% ------------------------------------------------------------
\newenvironment{gbpaginaespecial}
  {\gbclearrecto\thispagestyle{empty}\begin{center}}
  {\end{center}\gbclearrecto}

\newenvironment{gbpaginaimagen}
  {\gbclearrecto\thispagestyle{empty}%
   \begingroup\centering\setlength{\parindent}{0pt}}
  {\par\endgroup\gbclearrecto}

% ------------------------------------------------------------
% gbcolofon
% ------------------------------------------------------------
% VA EN VERSO. NO PUEDE SER UN \chapter*: EN LA CLASE book,
% \chapter FUERZA EL RECTO PORQUE @openright ESTÁ ACTIVO, Y
% CUALQUIER SALTO PREVIO QUEDA ANULADO.
% ------------------------------------------------------------
\newenvironment{gbcolofon}
  {\gbclearverso\thispagestyle{empty}\justifying}
  {\clearpage}

% ============================================================
% 5. ESTRUCTURA
% ============================================================

% ------------------------------------------------------------
% \gbContadoresSinCapitulo
% ------------------------------------------------------------
% SE LLAMA UNA VEZ AL ENTRAR EN PRELIMINARES Y UNA VEZ AL
% ENTRAR EN POSLIMINARES. NO POR PIEZA.
%
% EN ESAS DOS ZONAS \chapter NO EJECUTA \refstepcounter{chapter}
% (book.cls, \@chapter, RAMA \if@mainmatter). ESE COMANDO ES LO
% ÚNICO QUE INCREMENTA EL CONTADOR DE CAPÍTULO Y, A LA VEZ,
% REINICIA LOS QUE DEPENDEN DE ÉL: section, footnote, figure,
% table Y equation. SI NO CORRE, NO PASA NINGUNA DE LAS DOS.
%
% LLEVAR chapter A CERO ACTIVA LA GUARDA \ifnum \c@chapter>\z@
% DE \thefigure, \thetable Y \theequation, QUE OMITE EL PREFIJO.
% SIN ESTO, UNA FIGURA DE POSLIMINARES SALE "Figura 1.2", CON
% EL NÚMERO DEL ÚLTIMO CAPÍTULO DEL CUERPO. VERIFICADO.
%
% LAS FIGURAS Y TABLAS NO SE REINICIAN POR PIEZA A PROPÓSITO:
% EN PRELIMINARES LA NUMERACIÓN CORRIDA (1, 2, 3) ES LA CORRECTA
% Y REINICIAR PRODUCIRÍA DOS "Figura 1" CON \label EN COLISIÓN.
% LAS NOTAS AL PIE SÍ VAN POR PIEZA: \thefootnote ES ARÁBIGO
% PELADO Y NO HAY PREFIJO QUE COLISIONE.
% ------------------------------------------------------------
\newcommand{\gbContadoresSinCapitulo}{%
  \setcounter{chapter}{0}%
  \setcounter{section}{0}%
  \setcounter{figure}{0}%
  \setcounter{table}{0}%
  \setcounter{equation}{0}%
}

% ------------------------------------------------------------
% \gbPiezaSinNumerar
% ------------------------------------------------------------
% PARA UNA PIEZA SIN NUMERAR QUE VA EN LA ZONA DE CUERPO. EL
% CASO TÍPICO SON LAS CONCLUSIONES, QUE DEBEN PRECEDER A
% \appendix PARA QUE LOS APÉNDICES NUMEREN CON LETRA.
%
% COMO ESAS PIEZAS SE EMITEN CON \chapter*, NO CORRE
% \refstepcounter{chapter} Y NO SE REINICIA NADA: SU PRIMERA
% NOTA AL PIE ARRANCA DONDE LA DEJÓ EL CAPÍTULO ANTERIOR.
% VERIFICADO.
%
% NO SE TOCA EL CONTADOR DE CAPÍTULO: EL APÉNDICE POSTERIOR LO
% NECESITA INTACTO, Y \appendix YA LO LLEVA A CERO POR SU
% CUENTA. TAMPOCO SE TOCAN figure NI table: ESAS PIEZAS NO
% LLEVAN NI CUADROS NI FIGURAS.
% ------------------------------------------------------------
% TAMBIÉN APAGA LA NUMERACIÓN DE SECCIONES. SIN ESO, LA
% PRIMERA SECCIÓN DE LAS CONCLUSIONES SALE "1.1", QUE COLISIONA
% CON LA SECCIÓN 1.1 REAL DEL CAPÍTULO 1: \thesection ES
% \thechapter.\arabic{section} Y EL CONTADOR DE CAPÍTULO SIGUE
% EN 1. VERIFICADO EN EL SUMARIO.
% LA ZONA DE APÉNDICES LA VUELVE A ENCENDER.
\newcommand{\gbPiezaSinNumerar}{%
  \setcounter{footnote}{0}%
  \setcounter{section}{0}%
  \setcounter{secnumdepth}{0}%
}

% ------------------------------------------------------------
% \gbParte{título}{subtítulo}
% ------------------------------------------------------------
% LAS PARTES NO SON ARCHIVOS: SON METADATOS DE partes_libro Y
% LAS EMITE GAMBAS EN EL main.tex. NO LLEVAN TEXTO PROPIO.
% EL SUBTÍTULO PUEDE VENIR VACÍO. AL ÍNDICE VA SOLO EL TÍTULO.
% ------------------------------------------------------------
\newcommand{\gbParte}[2]{%
  \part[#1]{#1%
    \ifstrempty{#2}{}{\\[0.4em]{\normalsize\mdseries\itshape #2}}%
  }}

% ------------------------------------------------------------
% \gbCapitulo{título}{subtítulo}
% ------------------------------------------------------------
% EQUIVALENTE PARA CAPÍTULOS CON SUBTÍTULO.
% ------------------------------------------------------------
\newcommand{\gbCapitulo}[2]{%
  \chapter[#1]{#1%
    \ifstrempty{#2}{}{\\[0.3em]{\normalsize\mdseries\itshape #2}}%
  }}

% ============================================================
% 6. BLOQUES
% ============================================================

% CITA EN BLOQUE CON ATRIBUCIÓN OPCIONAL.
% NO ES UN EPÍGRAFE: DOCBOOK LOS DISTINGUE Y EL XSLT TAMBIÉN.
\newenvironment{gbcita}
  {\begin{quote}\small}
  {\end{quote}}
\newcommand{\gbatribucion}[1]{\par\vspace{2pt}\hfill\textup{\small#1}\par}

% EPÍGRAFE. LA ATRIBUCIÓN PUEDE SER UNA CITA BIBLIOGRÁFICA:
% EL XSLT RESUELVE UN <biblioref> DE ADENTRO A \textcite.
\newcommand{\gbepigrafe}[2]{%
  \par\vspace{0.5em}%
  \begin{flushright}%
  \begin{minipage}{0.72\linewidth}%
    \small\itshape #1\par
    \vspace{2pt}\raggedleft\upshape\footnotesize #2\par
  \end{minipage}%
  \end{flushright}%
  \vspace{0.8em}\par}

% RECUADRO LATERAL. UN ARGUMENTO: info | supplementary.
\newenvironment{gbsidebar}[1]
  {\begin{tcolorbox}[colback=black!4,colframe=black!4,
     boxrule=0pt,arc=2pt,left=6pt,right=6pt,top=4pt,bottom=4pt,
     before skip=6pt,after skip=6pt]\small\sffamily}
  {\end{tcolorbox}}

% VERSO. EL XSLT YA EMITE \\ ENTRE LÍNEAS.
\newenvironment{gbverso}
  {\begin{verse}\itshape}
  {\end{verse}}

% PARLAMENTO Y LOCUTOR. DOCBOOK 5.2 BASE NO TIENE dialogue NI
% drama: EL CANÓNICO USA para role="speech" CON
% emphasis role="speaker" (RC-DB-04).
\newenvironment{gbparlamento}
  {\par\begingroup\setlength{\parindent}{0pt}\hangindent=1.2em\hangafter=1}
  {\par\endgroup}
\newcommand{\gblocutor}[1]{\textsc{\MakeLowercase{#1}}\quad}

% ============================================================
% 7. LLAMADAS DE ATENCIÓN
% ============================================================
% CADA LLAMADA DE ATENCIÓN ES UN tcolorbox CON TÍTULO FIJO.
% NO SE USA UN .style DE pgfkeys PARAMETRIZADO: EL #1 DE UN
% \newcommand DENTRO DE UN .style LO CONSUME pgfkeys Y DA
% "You already have nine parameters", QUE NO DICE NADA SOBRE
% LA CAUSA. CON UN ENTORNO POR CLASE EL PROBLEMA NO EXISTE.
\tcbset{gbadmon/.style={
  colback=black!3, colframe=black!25,
  boxrule=0.4pt, arc=1pt,
  left=6pt, right=6pt, top=4pt, bottom=4pt,
  before skip=6pt, after skip=6pt,
  fonttitle=\bfseries\sffamily\small}}

\newenvironment{gbnota}
  {\begin{tcolorbox}[gbadmon,title={Nota}]\small}{\end{tcolorbox}}
\newenvironment{gbconsejo}
  {\begin{tcolorbox}[gbadmon,title={Sugerencia}]\small}{\end{tcolorbox}}
\newenvironment{gbadvertencia}
  {\begin{tcolorbox}[gbadmon,title={Advertencia}]\small}{\end{tcolorbox}}
\newenvironment{gbimportante}
  {\begin{tcolorbox}[gbadmon,title={Importante}]\small}{\end{tcolorbox}}
\newenvironment{gbprecaucion}
  {\begin{tcolorbox}[gbadmon,title={Precaución}]\small}{\end{tcolorbox}}

% ============================================================
% 8. CÓDIGO
% ------------------------------------------------------------
% listings Y NO minted: minted EXIGE --shell-escape Y Pygments,
% Y ESO ROMPE LA AUTONOMÍA DEL PROYECTO EN OTRA MÁQUINA.
% ============================================================
\lstset{
  basicstyle=\ttfamily\footnotesize,
  breaklines=true,
  columns=fullflexible,
  frame=none,
  showstringspaces=false,
  extendedchars=true,
  inputencoding=utf8
}

% ============================================================
% 9. COLOFÓN
% ------------------------------------------------------------
% LAS FUENTES Y LA VERSIÓN SON VALORES, NO TEXTO ESCRITO. SI
% UNA EDITORIAL REEMPLAZA LA ESTÉTICA Y CAMBIA LAS FUENTES,
% EL COLOFÓN SIGUE DICIENDO LA VERDAD SIN QUE NADIE SE
% ACUERDE DE EDITARLO.
% LA FECHA ES LA DE PUBLICACIÓN Y NO LA DEL DÍA: UNA MARCA DE
% TIEMPO HARÍA QUE CADA COMPILACIÓN PRODUJERA UN DIFF AUNQUE
% NADA HUBIERA CAMBIADO.
% ============================================================
% SI LA ESTÉTICA NO LAS DECLARA, EL COLOFÓN LO DICE EN VEZ DE DAR UN
% "Undefined control sequence" QUE NO EXPLICA NADA.
\providecommand{\gbfuentetexto}{(sin declarar)}
\providecommand{\gbfuentetitulos}{(sin declarar)}
\providecommand{\gbversion}{(sin version)}

% EL RECUADRO SE DEFINE COMO ESTILO Y NO EN LÍNEA: ASÍ UNA EDITORIAL
% QUE QUIERA OTRO ASPECTO REDEFINE gbcolofoncaja SIN TOCAR EL TEXTO,
% QUE ES LO QUE EL CONTRATO GARANTIZA.
\tcbset{gbcolofoncaja/.style={
  colback=black!7,
  colframe=black!50,
  boxrule=0.7pt,
  arc=3mm,
  left=8pt, right=8pt, top=8pt, bottom=8pt,
  boxsep=0pt}}

\newcommand{\gbTextoColofon}{%
  \begin{tcolorbox}[gbcolofoncaja]
    La composición tipográfica de este libro se realizó con
    gbpublisher versión~\gbversion, bajo un modelo de fuente única que
    garantiza la trazabilidad entre el original y cada una de sus
    salidas (HTML, ePub y PDF). Las familias tipográficas utilizadas son
    \gbfuentetexto\ para el texto e \gbfuentetitulos\ para los
    títulos. El proyecto se encuentra disponible en
    \url{https://github.com/albertomoyano/gbpublisher}.
  \end{tcolorbox}}

   entornos y macros del XSLT
%   % ============================================================
% preambulo-estetica.tex
% ------------------------------------------------------------
% CAPA REEMPLAZABLE. ES LA QUE UNA EDITORIAL CAMBIA POR LA SUYA
% CUANDO QUIERE MANTENER SU IDENTIDAD VISUAL. NO DEBE DEFINIR
% NINGUNO DE LOS ENTORNOS DEL CONTRATO.
% ------------------------------------------------------------
% ESTE ARCHIVO ES EL ESQUELETO. LOS BLOQUES DE DISEÑO QUE YA
% EXISTEN EN EL preambulo.tex EN USO —titlesec, titletoc,
% estilos de página, diseño de notas al pie, newfloat,
% imakeidx, glossaries, froufrou, siunitx— SE PEGAN EN EL
% APARTADO 8. NO ESTÁN ACÁ PORQUE SON DECISIONES DE DISEÑO
% PROPIAS Y NO CORRESPONDE REESCRIBIRLAS.
% ============================================================
% LAS CUATRO PRESTACIONES
% ------------------------------------------------------------
% El derivado emite PRESTACIONES y no un «modo». Así, agregar
% un estado editorial nuevo no obliga a tocar esta capa, y cada
% bandera dice qué hace en vez de cómo se llama.
%
%   estado_libro     showframe  homeoarchy  enlaces  cruces
%   borrador             si         si         si      no
%   en_correccion        no         si         si      no
%   en_produccion        no         no         si      no
%   publicado            no         no         no      si
%
% LAS CUATRO TOCAN OPCIONES DE PAQUETE, QUE NO SE PUEDEN
% SUPERPONER DESPUES DE CARGADO EL PAQUETE. POR ESO SON
% BANDERAS DEL DERIVADO Y NO REEMPLAZOS DE TEXTO.
% LA EXCEPCION ES hyperref: SE CARGA UNA SOLA VEZ EN EL
% CONTRATO Y LOS COLORES VAN POR \hypersetup, QUE SI SE
% SUPERPONE.
% ============================================================
% CORRECCIONES APLICADAS RESPECTO DEL PREAMBULO EN USO
%   - FUERA \usepackage[utf8]{luainputenc}: LuaLaTeX CONSUME
%     UTF-8 NATIVAMENTE (SC-02).
%   - FUERA \usepackage{svg}: EXIGE Inkscape Y --shell-escape,
%     Y EN LIBROS SOLO SE USAN PNG, JPG Y PDF.
%   - \renewcommand{\footnoterule} ESTABA DEFINIDO DOS VECES
%     CON EL MISMO CUERPO. QUEDA UNO.
%   - babel LO CARGA EL CONTRATO, ANTES DE csquotes Y biblatex,
%     QUE ES EL ORDEN QUE PIDE LA DOCUMENTACION.
%   - luatexja SOLO SI EL LIBRO LO NECESITA.
% ============================================================

\usepackage{ifthen}
\usepackage{calc}

% ============================================================
% 1. GEOMETRIA
% ------------------------------------------------------------
% Los valores vienen del derivado.
%
% showframe ES OPCION DE PAQUETE Y NO SE PUEDE SUPERPONER: SE
% PASA CON \PassOptionsToPackage PARA NO DUPLICAR LA LISTA
% ENTERA DE OPCIONES EN DOS RAMAS DE UN CONDICIONAL.
%
% Hoy los cuatro margenes son iguales en los 45 formatos de la
% base, asi que inner/outer da lo mismo que centering; se usan
% igual para que el dia que se carguen asimetricos no haya que
% tocar nada. En un libro encuadernado el margen interior suele
% necesitar mas aire que el exterior.
% ============================================================
\ifgbshowframe
  \PassOptionsToPackage{showframe}{geometry}
\fi

\usepackage[
  paperwidth=\gbanchopagina,
  paperheight=\gbaltopagina,
  inner=\gbmargeninterior,
  outer=\gbmargenexterior,
  top=\gbmargensuperior,
  bottom=\gbmargeninferior,
  includehead,
  includefoot,
  headsep=14pt,
  footskip=24pt
]{geometry}

% ------------------------------------------------------------
% CRUCES DE CORTE
% ------------------------------------------------------------
% La hoja de montaje es mayor que el formato final: el derivado
% la calcula sumando el margen de cruces por lado. Las cruces
% solo necesitan lugar, y ese lugar es el mismo para cualquier
% formato.
%
% crop DECLARA LA HOJA DE MONTAJE Y geometry EL FORMATO FINAL:
% SON DOS PARES DE MEDIDAS INDEPENDIENTES, NO UNA DERIVADA DE
% LA OTRA. VERIFICADO: crop EXPANDE MACROS EN SUS OPCIONES
% keyval, ASI QUE NO HACE FALTA \crop[...] APARTE.
% ------------------------------------------------------------
\ifgbcruces
  \usepackage[width=\gbanchomontaje,height=\gbaltomontaje,cam,center]{crop}
\fi

% ============================================================
% 2. IDIOMA
% ------------------------------------------------------------
% babel lo carga el contrato, antes de csquotes y biblatex.
% Aca solo se cambian las cadenas, que si son diseno.
% ============================================================
\addto\captionsspanish{\renewcommand{\contentsname}{Sumario}}

% ============================================================
% 3. TIPOGRAFIA
% ------------------------------------------------------------
% Las fuentes son de la estetica, no de la base. Por eso cuando
% una editorial reemplaza esta capa, cambia las fuentes y el
% colofon las nombra correctamente sin que haya que tocar nada
% mas: el contrato solo las lee.
% ============================================================
\usepackage{fontspec}
\usepackage[final,babel=true]{microtype}

\def\gbfuentetexto{Libertinus Serif}
\def\gbfuentetitulos{IBM Plex Sans Condensed}
\def\gbfuentemono{IBM Plex Mono}

\setmainfont{\gbfuentetexto}[
  Numbers={OldStyle,Proportional},
  Ligatures={TeX,Common,Discretionary},Scale=1.18]
\setsansfont{\gbfuentetitulos}[
Scale=MatchLowercase,
Ligatures=TeX,
Extension=.otf,
UprightFont=*-Regular,
ItalicFont=*-Italic,
BoldFont=*-SemiBold,
BoldItalicFont=*-SemiBoldItalic]
\setmonofont{\gbfuentemono}[
Scale=0.91,
Extension=.otf,
UprightFont=*-Regular,
ItalicFont = IBMPlexMono-Italic.otf,
BoldFont = IBMPlexMono-Bold.otf,
BoldItalicFont = IBMPlexMono-BoldItalic.otf]

%\setsansfont{\gbfuentetitulos}[Scale=MatchLowercase,Ligatures=TeX]
%\setmonofont{\gbfuentemono}[Scale=0.91]

\usepackage{unicode-math}
\setmathfont{Libertinus Math}[Scale=MatchLowercase]

% SOPORTE CJK SOLO SI EL LIBRO LO NECESITA: luatexja PESA EN
% CADA COMPILACION Y CASI NINGUN LIBRO LO USA.
\ifgbcjk
  \usepackage{luatexja-fontspec}
  \setmainjfont{FandolSong}
\fi

\renewcommand{\normalsize}{\fontsize{10pt}{14pt}\selectfont}
\normalsize

% ============================================================
% 4. LISTAS, NOTAS Y FILETES
% ============================================================
% REEMPLAZA EL APARTADO 4 DE preambulo-estetica.tex
%
% QUÉ CAMBIÓ RESPECTO DEL ESQUELETO, Y POR QUÉ
% ------------------------------------------------------------
% 1. SE SACÓ hang DE footmisc
%
%    hang SANGRA EL TEXTO DE LA NOTA HACIA ADENTRO PARA HACERLE
%    LUGAR A LA MARCA. LO QUE SE BUSCA ES LO CONTRARIO: LA MARCA
%    COLGADA EN EL MARGEN Y EL TEXTO OCUPANDO EL ANCHO COMPLETO
%    DE LA CAJA. ESO LO PRODUCE \deffootnote DE scrextend.
%
%    hang Y \deffootnote GOBIERNAN LA MISMA SANGRÍA POR
%    MECANISMOS DISTINTOS, ASÍ QUE CARGAR LOS DOS DEJA EL
%    RESULTADO A MERCED DE CUÁL SE APLIQUE ÚLTIMO. QUEDA UNO.
%
%    bottom Y stable SÍ SE CONSERVAN: FIJAN LAS NOTAS AL PIE DE
%    LA CAJA Y LAS PERMITEN EN LOS TÍTULOS, Y NINGUNO TOCA LA
%    SANGRÍA.
%
% 2. SE SACARON LOS DOS \patchcmd SOBRE \@footnotetext
%
%    VERIFICADO: LOS DOS FALLAN. scrextend REDEFINE
%    \@footnotetext, ASÍ QUE CUANDO EL PARCHE BUSCA \footnotesize
%    O \@MM ESOS TOKENS YA NO ESTÁN EN LA DEFINICIÓN.
%
%    EL PRIMERO SE REEMPLAZA POR \setkomafont{footnote}, QUE ES
%    EL MECANISMO PROPIO DE KOMA Y NO DEPENDE DE PARCHEAR NADA.
%    VERIFICADO: CON \small LA MISMA NOTA OCUPA DOS PÁGINAS Y CON
%    \footnotesize UNA.
%
%    EL SEGUNDO SE REEMPLAZA POR \interfootnotelinepenalty, EL
%    PARÁMETRO DEL KERNEL QUE GOBIERNA ESE CORTE. ADVERTENCIA
%    HONESTA: NO PUDE DEMOSTRAR QUE CAMBIE NADA. UNA NOTA DE
%    CATORCE PÁRRAFOS SE PARTIÓ IGUAL CON EL PARÁMETRO Y SIN ÉL.
%    QUEDA PORQUE ES EL MECANISMO DOCUMENTADO Y NO HACE DAÑO,
%    PERO SI EL COMPORTAMIENTO CON NOTAS LARGAS NO ES EL DEL
%    PREÁMBULO LEGACY, ESTE ES EL LUGAR A REVISAR.
% ============================================================

\usepackage{enumitem}
\setlist{nosep,topsep=4pt}

% ------------------------------------------------------------
% NOTAS AL PIE
% ------------------------------------------------------------
\usepackage[bottom,stable]{footmisc}
\usepackage{scrextend}

% SEPARACIÓN ENTRE NOTAS
\setlength{\footnotesep}{10pt}

% FILETE SEPARADOR: 0,5 pt DE GROSOR Y 40% DEL ANCHO DE COLUMNA
\renewcommand{\footnoterule}{%
	\kern-3pt
	\hrule height 0.5pt width 0.4\columnwidth
	\kern6pt}

% NUMERACIÓN: ENTRE CORCHETES, PALO SECO, CUERPO scriptsize
\renewcommand*{\thefootnote}{\scriptsize\sffamily{[\arabic{footnote}]}}

% CUERPO DE LA NOTA
\setkomafont{footnote}{\small}

% ------------------------------------------------------------
% MARCA COLGADA EN EL MARGEN
% ------------------------------------------------------------
% \llap SACA LA MARCA FUERA DE LA CAJA HACIA LA IZQUIERDA, ASÍ EL
% TEXTO DE LA NOTA EMPIEZA EN EL MARGEN Y OCUPA EL ANCHO COMPLETO.
%
% LOS TRES ARGUMENTOS DE \deffootnote SON: SANGRÍA DE LA PRIMERA
% LÍNEA, SANGRÍA DE LAS SIGUIENTES, Y CÓMO SE COMPONE LA MARCA.
% ------------------------------------------------------------
\newcommand*\gbmarcanotaespacio{0em}

\deffootnote{\gbmarcanotaespacio}
{\parindent}
{%
	\makebox[\gbmarcanotaespacio][r]{\llap{\thefootnotemark\quad}}%
}

% ------------------------------------------------------------
% NOTAS LARGAS PARTIDAS ENTRE PÁGINAS
% ------------------------------------------------------------
% Penalización de partir una nota entre dos páginas. Por omisión
% LaTeX pone 100 en la clase book, y el valor alto que usa el
% kernel en otros contextos hace que una nota muy larga desplace
% texto en lugar de partirse.
%
% VER LA ADVERTENCIA DEL ENCABEZADO: NO PUDE MEDIR DIFERENCIA.
% ------------------------------------------------------------
\interfootnotelinepenalty=100

% ------------------------------------------------------------
% RESTO DEL APARTADO
% ------------------------------------------------------------
\usepackage[labelfont=bf,font=small,labelsep=period,format=plain]{caption}
\usepackage{needspace}

\setcounter{tocdepth}{1}
\setcounter{secnumdepth}{3}
\renewcommand{\thepart}{\arabic{part}}

\raggedbottom
\clubpenalty=10000
\widowpenalty=10000

% ============================================================
% 5. MARCAS TIPOGRAFICAS DE TRABAJO
% ------------------------------------------------------------
% impnattypo marca homeoarquias y problemas de principio y fin
% de linea. La opcion draft es opcion de paquete y no se puede
% superponer: por eso la carga va condicionada y no comentada
% a mano.
% ============================================================
\ifgbhomeoarchy
  \IfFileExists{impnattypo.sty}{%
    \usepackage[hyphenation,homeoarchy,homeoarchywordcolor=orange,
                homeoarchycharcolor=orange,draft]{impnattypo}}{%
    \PackageWarning{gbpublisher}{impnattypo no instalado:
      se pierden las marcas tipograficas de trabajo}}
  \overfullrule=5pt
\fi

% ============================================================
% 6. COLOR DE LOS ENLACES
% ------------------------------------------------------------
% Solo el color. hyperref lo carga preambulo-hyperref.tex, que
% va después de esta capa y lee \ifgbenlaces para decidir entre
% enlaces coloreados y hidelinks.
% ============================================================
\usepackage{xcolor}
\definecolor{gbenlace}{RGB}{160,0,120}

% ============================================================
% 7. CONSTANCIA EN EL LOG
% ------------------------------------------------------------
% Deja escrito en el .log con que prestaciones se compuso.
% Cuando un PDF sale distinto de lo esperado, esta linea dice
% por que sin tener que abrir el derivado.
% ============================================================
\typeout{[gbpublisher] prestaciones:
  showframe=\ifgbshowframe si\else no\fi,
  homeoarchy=\ifgbhomeoarchy si\else no\fi,
  enlaces=\ifgbenlaces si\else no\fi,
  cruces=\ifgbcruces si\else no\fi}

% ============================================================
% 8. AQUI VAN LOS BLOQUES DE DISENO DEL PREAMBULO EN USO
% ------------------------------------------------------------
%   - titlesec y los estilos de titulo
%   - titletoc: partes, capitulos y secciones del sumario
%   - titleps: estilos de pagina y folios corridos
%     (\gbtituloabreviado ES LO QUE CORRESPONDE EN EL FOLIO
%      CORRIDO, Y HOY EL PREAMBULO NO LO CONSUME)
%   - newfloat: el entorno flotante "imagen"
%   - imakeidx + xindy: names (lo llena esindex desde biblatex),
%     onomastico y concepto (los carga el editor)
%   - glossaries: siglas y glosario
%   - siunitx, froufrou, qrcode, pdfpages, rotating
% ============================================================

% Cita con cambio de márgenes y tamaño tipografico
\renewenvironment{quote}
{\normalsize\list{}{\sf\leftmargin=14pt \rightmargin=0pt}%
	\item\relax}
{\endlist}

% Configuración de los índices
\usepackage[xindy]{imakeidx}
\makeindex[name=names, title={}, intoc=false]
\makeindex[name=onomastico, title={}, intoc=false]
\makeindex[name=concepto, title={}, intoc=false]


% ============================================================
% PERSONALIZACIÓN DE LOS ÍNDICES DE PARTES, CAPÍTULOS Y SECCIONES
% ------------------------------------------------------------
% Paquete requerido:
\usepackage{titletoc}
%
% OBJETIVO GENERAL:
% - Controlar la apariencia de los índices (Tabla de Contenidos)
%   para partes, capítulos y secciones.
% - Ajustar sangrías, tipografía, alineación y estilo de los
%   números de página.
% ============================================================

% ------------------------------------------------------------
% PARTES
% ------------------------------------------------------------
% \titlecontents{part}[<left>]{<above-code>}{<numbered-entry-format>}
%                        {<numberless-entry-format>}{<filler-page>}
%
% Argumentos usados:
%   [0em]                    -> Sangría izquierda de la entrada.
%   {\addvspace{5pt}\sf\bfseries\normalsize\selectfont\filright}
%                             -> Espacio vertical antes de la entrada (5pt),
%                                fuente sans serif, negrita, tamaño normal,
%                                alineación a la derecha (\filright).
%   {\contentslabel[\thecontentslabel]{2.5pc}}
%                             -> Formato de la etiqueta numérica (ej. Parte I)
%                                con ancho de 2.5pc.
%   {}                        -> Formato de entrada sin número.
%   {}                        -> No hay reglas ni página aquí.
\titlecontents{part}[0em]
{\addvspace{5pt}\sf\bfseries\normalsize\selectfont\filright}
{\contentslabel[\thecontentslabel]{2.5pc}}
{}
{}

% ------------------------------------------------------------
% CAPÍTULOS
% ------------------------------------------------------------
% Formato general:
%   - Sangría: 1.5pc
%   - Espacio vertical: 0.4em antes de cada capítulo
%   - Fuente: sans serif, alineación a la derecha (\filright)
%   - Número del capítulo alineado con margen de 1.5pc
%   - Separador de puntos hasta el número de página (\titlerule)
%
% Argumentos:
%   [1.5pc] -> sangría izquierda del capítulo
%   {\addvspace{.4em}\sf\selectfont\filright} -> formato general de la entrada
%   {\contentslabel{1.5pc}} -> ancho reservado para número
%   {\hspace*{-1.5pc}} -> corrección si no tiene número
%   {\titlerule*[1pc]{.}\contentspage[\hspace*{-4pc} {\rm\small\thecontentspage}]}
%      -> regla de puntos de relleno hasta el número de página,
%         página en tamaño pequeño y ligeramente desplazada
\titlecontents{chapter}[1.5pc]
{\addvspace{.4em}\sf\selectfont\filright}
{\contentslabel{1.5pc}}
{\hspace*{-1.5pc}}%
{\titlerule*[1pc]{.}\contentspage[\hspace*{-4pc} {\rm\small\thecontentspage}]}%
[]

% ------------------------------------------------------------
% SECCIONES
% ------------------------------------------------------------
% Formato general:
%   - Sangría: 4.5pc
%   - Fuente: \small, alineación a la derecha
%   - Número de sección: 2.5pc
%   - Ajuste de sangría si no hay número: -2.5pc
%   - Separador de puntos hasta el número de página (\titlerule)
\titlecontents{section}[4.5pc]
{\small\filright}                  % Formato general
{\contentslabel{2.5pc}}            % Número con ancho 2.5pc
{\hspace*{-2.5pc}}                 % Si no tiene número
{\titlerule*[1pc]{.}\contentspage} % Regla de puntos hasta la página


% ============================================================
% CONFIGURACIÓN DE ESTILOS DE TÍTULOS CON titlesec
% ------------------------------------------------------------
% Paquete requerido:
\usepackage[sf,bf,compact,topmarks,calcwidth,pagestyles,clearempty,newlinetospace]{titlesec}
%
% OBJETIVO GENERAL:
% - Personalizar la apariencia de partes, capítulos, secciones,
%   subsecciones, subsubsecciones, párrafos y subpárrafos.
% - Controlar fuente, tamaño, alineación, sangrías y espaciado.
% ============================================================

% ============================================================
% DISEÑO DE PARTES
% ------------------------------------------------------------
% Redefinición de \@part y \@spart para controlar:
% - Numeración de partes (si corresponde)
% - Entrada en el TOC
% - Centrado del título y subtítulo de la parte
% - Tamaño de fuente y estilo tipográfico
%
% \@part[#1]#2:
%   - #1: título corto para TOC y encabezado
%   - #2: título completo para mostrar en la página
%   - Centrado, fuente sans serif (\sf), \LARGE para nombre de parte
%     y \Large para título completo
%
% \@spart#1:
%   - Para partes no numeradas
%   - Solo muestra el título centrado, \LARGE, sans serif
%
\makeatletter
\def\@part[#1]#2{%
	\ifnum \c@secnumdepth >-2\relax
	\refstepcounter{part}%
	\addcontentsline{toc}{part}{Parte \thepart\hspace{1em}#1}%
	\else
	\addcontentsline{toc}{part}{#1}%
	\fi
	\markboth{}{}%
	{\centering
		\interlinepenalty \@M
		\ifnum \c@secnumdepth >-2\relax
		\sf\LARGE\selectfont \partname~\thepart
		\par\vskip 20\p@%
		\fi
		\sf\Large\selectfont
		#2\par}%
	\@endpart
}

\def\@spart#1{%
	\markboth{}{}%
	{\centering
		\interlinepenalty \@M
		\sf\LARGE\selectfont
		#1\par}%
	\@endpart
}
\makeatother

% ============================================================
% DISEÑO DE CAPÍTULO
% ------------------------------------------------------------
% \titleformat{\chapter}[display]{...}{...}{sep}{before-code}[after-code]
% - [display]: título en bloque separado (no inline)
% - \Large: tamaño base del título
% - \vspace{-2.5cm}: ajusta altura vertical (sube el título)
% - \centering: centrado
% - \textsc{\MakeLowercase\chaptertitlename}: 'Capítulo' en versalitas
% - \thechapter: número del capítulo
% - 1.5cm: separación entre número y texto
% - \filright\sf\LARGE: alineación a la derecha, fuente sans serif, tamaño LARGER
\titleformat{\chapter}[display]
{\Large}
{\centering{\textsc{\MakeLowercase\chaptertitlename}}~\thechapter}
{1.5cm}
{\filright\sf\LARGE}
[]
% ------------------------------------------------------------
% ESPACIADO DEL TÍTULO DE CAPÍTULO
% ------------------------------------------------------------
% \titlespacing*{\chapter}{izquierda}{antes}{después}
%
% ANTES = 0pt HACE QUE «capítulo N» ARRANQUE EN EL BORDE SUPERIOR
% DE LA CAJA. LA CLASE book RESERVA 50pt POR OMISIÓN.
%
% LA VERSIÓN CON ASTERISCO SUPRIME LA SANGRÍA DEL PRIMER PÁRRAFO
% DESPUÉS DEL TÍTULO.
%
% NO SE USA \vspace NEGATIVO DENTRO DEL \titleformat: ESO EMPUJA
% HACIA ARRIBA DESPUÉS DE QUE LaTeX YA RESERVÓ EL ESPACIO, Y EL
% RESULTADO DEPENDE DEL CONTEXTO. titlesec TIENE EL PARÁMETRO
% PARA DECIRLO DIRECTO.
% ------------------------------------------------------------
\titlespacing*{\chapter}{0pt}{0pt}{30pt}

% ============================================================
% DISEÑO DE SECCIÓN
% ------------------------------------------------------------
% Fuente: sans serif, negrita, tamaño grande (\large), alineación a la derecha
\titleformat{\section}[hang]
{\sf\bfseries\large\raggedright}
{\thesection}{.5em}{}[]
\titlespacing{\section}
{\parindent}{18pt plus 0.5pt minus 0.5pt}{6.75pt}

% ============================================================
% DISEÑO DE SUBSECCIÓN
% ------------------------------------------------------------
\titleformat{\subsection}[hang]
{\sf\large\raggedright}
{\thesubsection}{.5em}{}[]
\titlespacing{\subsection}
{\parindent}{18pt plus 0.5pt minus 0.5pt}{6.75pt}

% ============================================================
% DISEÑO DE SUBSUBSECCIÓN
% ------------------------------------------------------------
\titleformat{\subsubsection}[hang]
{\rm\bfseries\normalsize\raggedright}
{\thesubsubsection}{.5em}{}[]
\titlespacing{\subsubsection}
{\parindent}{18pt plus 0.5pt minus 0.5pt}{6.75pt}

% ============================================================
% DISEÑO DE PÁRRAFO
% ------------------------------------------------------------
% \paragraph: run-in (alineado inline, no bloque)
% - Negrita
% - Sin número
% - Raya al final del título (---)
\titlespacing{\paragraph}
{0pt}                % margen izquierda
{6.75pt plus 0.5pt minus 0.5pt} % espacio antes
{0pt}                % espacio después
\titleformat{\paragraph}[runin]
{\bfseries}{}{0em}{}[\mbox{ --- }]

% ============================================================
% DISEÑO DE SUBPÁRRAFO
% ------------------------------------------------------------
% \subparagraph: estilo colapsado, centrado
% - Fuente sans serif, negrita, tamaño \large
\titlespacing{\subparagraph}
{\parindent}{18pt plus 0.5pt minus 0.5pt}{6.75pt}
\titleformat{\subparagraph}[hang]
{\sf\bfseries\large\centering}
{\thesubparagraph}{.5em}{}[]

   diseño, reemplazable
%   % ============================================================
% preambulo-hyperref.tex
% ------------------------------------------------------------
% CUARTA Y ÚLTIMA CAPA DEL PREÁMBULO. VA DESPUÉS DE LA ESTÉTICA.
%
%   % ============================================================
% preambulo-proyecto.tex
% ------------------------------------------------------------
% ARCHIVO DERIVADO — GENERADO POR gbpublisher
% NO EDITAR A MANO: SE REESCRIBE EN CADA GENERACIÓN DEL
% main.tex. LOS CAMBIOS SE HACEN EN LA BASE.
%
% SÍ SE VERSIONA: EL COMMIT ES LA GARANTÍA DE TRAZA DE QUÉ SE
% PUBLICÓ. POR ESO NO LLEVA MARCA DE TIEMPO.
% ------------------------------------------------------------
% ESTA CAPA SOLO DEFINE VALORES. NO CARGA NINGÚN PAQUETE.
% VA PRIMERA PORQUE geometry Y fontspec NECESITAN LOS VALORES
% EN EL MOMENTO DEL \usepackage.
% ============================================================

% --- PRESTACIONES DE SALIDA (derivadas de estado_libro) ---
\newif\ifgbshowframe    \gbshowframefalse
\newif\ifgbhomeoarchy   \gbhomeoarchytrue
\newif\ifgbenlaces      \gbenlacestrue
\newif\ifgbcruces       \gbcrucesfalse

% --- FORMATO DE PÁGINA (formatos_pdf id=32) ---
\def\gbanchopagina{15.50cm}
\def\gbaltopagina{23.00cm}
\def\gbmargeninterior{2.00cm}
\def\gbmargenexterior{2.00cm}
\def\gbmargensuperior{2.00cm}
\def\gbmargeninferior{2.00cm}
\def\gbnombreformato{Estudio 2A / 04}
\def\gbanchomontaje{18.00truecm}
\def\gbaltomontaje{25.50truecm}

% --- IDIOMA ---
\def\gbidiomas{main=spanish}
\def\gbidiomaprincipal{spanish}
\newif\ifgbcjk   \gbcjkfalse

% --- BIBLIOGRAFÍA ---
\def\estandar{apa}
\newif\ifgbbibporcapitulo   \gbbibporcapitulotrue
\def\gbarchivobib{./referencias/ref-l-DUEK.bib}

% --- METADATOS DEL PDF ---
\def\gbtitulo{Un tiempo diferente}
\def\gbsubtitulo{Vida cotidiana y pandemia}
\def\gbtituloabreviado{}
\def\gbautores{Moguillansky, Marina (comp.); Duek, Carolina (comp.)}
\def\gbeditorial{Ediciones Imago Mundi}
\def\gbciudad{Buenos Aires}
\def\gbanio{2026}
\def\gbisbn{978-950-793-487-2}
\def\gbisbnelectronico{}
\def\gbdoi{}
\def\gblicencia{}
\def\gburllicencia{}
\def\gbpalabrasclave{}

% --- VERSIÓN DE gbpublisher ---
\def\gbversion{0.0.332}
   valores de la base
%   % ============================================================
% preambulo-contrato.tex
% ------------------------------------------------------------
% CONTRATO VERSIÓN 1
% ------------------------------------------------------------
% DEFINE TODO LO QUE docbook-to-latex.xsl PUEDE EMITIR EN LOS
% FRAGMENTOS DE CAPÍTULO. ES LA CAPA QUE UNA EDITORIAL
% CONSERVA CUANDO REEMPLAZA LA ESTÉTICA: SI FALTA UNA DE
% ESTAS DEFINICIONES, EL LIBRO NO COMPILA; SI ESTÁ DEFINIDA
% DISTINTO, COMPILA Y SALE MAL.
%
% CUANDO SE AGREGUE VOCABULARIO NUEVO, SUBIR LA VERSIÓN DEL
% CONTRATO. EL main.tex DECLARA LA QUE NECESITA Y ESTA CAPA
% LA VERIFICA: ASÍ UN PREÁMBULO VIEJO FRENA EN LA PRIMERA
% LÍNEA CON UN MENSAJE ÚTIL, EN VEZ DE DAR UN
% "Undefined control sequence" EN LA PÁGINA 140.
% ------------------------------------------------------------
% DEPENDE DE preambulo-proyecto.tex, QUE VA ANTES.
% ============================================================

\def\gbContratoVersion{1}

% ============================================================
% \gbRequiereContrato{N}
% PROPÓSITO : EL main.tex DECLARA QUÉ VERSIÓN NECESITA.
% PARÁMETRO : N — versión requerida
% ============================================================
\newcommand{\gbRequiereContrato}[1]{%
  \ifnum#1>\gbContratoVersion\relax
    \PackageError{gbpublisher}{%
      Este preambulo implementa el contrato \gbContratoVersion\space
      y el libro requiere el #1}{%
      Actualizar preambulo-contrato.tex.}%
  \fi}

% ============================================================
% HOJAS DE CORTESÍA
% ------------------------------------------------------------
% Dos páginas en blanco al inicio —una hoja— que todo libro
% lleva. Entran en la paginación: son i y ii, y el sumario
% empieza en iii.
%
% \hbox{} ES IMPRESCINDIBLE: UNA PÁGINA SIN NINGÚN MATERIAL NO
% SE EMITE, ASÍ QUE UN \newpage SOBRE UNA PÁGINA VACÍA NO
% PRODUCE NADA Y EL CONTADOR NO AVANZA. LA CAJA VACÍA DA EL
% MÍNIMO DE CONTENIDO QUE OBLIGA A COMPONERLA.
%
% \thispagestyle{empty} LES QUITA EL FOLIO: SE CUENTAN PERO NO
% SE IMPRIME EL NÚMERO.
% ============================================================
\newcommand{\gbHojasCortesia}{%
	\thispagestyle{empty}\hbox{}\newpage
	\thispagestyle{empty}\hbox{}\newpage}

% ============================================================
% ENCABEZADO DE PIEZA GENERADA
% ------------------------------------------------------------
% Para las piezas sin cuerpo que llevan título propio. El título
% sale de capitulos.titulo_capitulo, que es donde el editor lo
% escribe, y no del preámbulo.
%
% \markboth actualiza el folio corrido, que \chapter* no toca.
% ============================================================
\newcommand{\gbEncabezadoPieza}[1]{%
	\chapter*{#1}%
	\addcontentsline{toc}{chapter}{#1}%
	\markboth{\MakeUppercase{#1}}{\MakeUppercase{#1}}}

% ============================================================
% ÍNDICE CON SU ENCABEZADO
% ------------------------------------------------------------
% \printindex HACE UN \clearpage INTERNO AUNQUE EL ÍNDICE SE
% DECLARE CON title={}. SIN NEUTRALIZARLO, EL ENCABEZADO QUEDA
% SOLO EN UNA PÁGINA Y LA LISTA EMPIEZA EN LA SIGUIENTE.
%
% EL \let VA DENTRO DE UN GRUPO: FUERA DE ÉL \clearpage TIENE QUE
% SEGUIR FUNCIONANDO, Y EL COLOFÓN QUE VIENE DESPUÉS LO NECESITA.
% ============================================================
\newcommand{\gbIndice}[2]{%
	\gbEncabezadoPieza{#1}%
	\begingroup
	\let\clearpage\relax
	\raggedright
	\small
	\printindex[#2]%
	\endgroup}

% ============================================================
% 1. PAQUETES QUE EL CONTRATO NECESITA
% ------------------------------------------------------------
% SOLO LO QUE HACE FALTA PARA QUE UN FRAGMENTO COMPILE. LA
% ESTÉTICA CARGA LO SUYO APARTE. TODOS ESTÁN EN TeX Live Y
% NINGUNO PIDE --shell-escape.
% ============================================================
\usepackage{etoolbox}
\usepackage{graphicx}
\usepackage{booktabs}
\usepackage{array}
\usepackage{listings}
\usepackage{ragged2e}
\usepackage{xcolor}
\usepackage[most]{tcolorbox}

% ------------------------------------------------------------
% IDIOMA
% ------------------------------------------------------------
% babel VA EN EL CONTRATO Y NO EN LA ESTÉTICA, POR DOS RAZONES.
% LA PRIMERA ES DE ORDEN: biblatex PIDE CARGARSE DESPUÉS DE
% babel Y csquotes, Y LA ESTÉTICA VA DESPUÉS DE ESTA CAPA.
% LA SEGUNDA ES DE FONDO: EL IDIOMA DECIDE LAS COMILLAS QUE
% PRODUCE \enquote, LA PARTICIÓN DE PALABRAS Y LAS CADENAS DE
% LA BIBLIOGRAFÍA. SON GARANTÍAS DEL CONTRATO, NO DECISIONES
% DE DISEÑO: UNA EDITORIAL QUE REEMPLAZA LA ESTÉTICA NO DEBE
% PODER CAMBIAR EL IDIOMA DEL LIBRO.
\usepackage[\gbidiomas]{babel}
\frenchspacing

% COMILLAS SEGÚN EL IDIOMA ACTIVO: EN ESPAÑOL DA « ».
% EL XSLT EMITE \enquote Y NUNCA COMILLAS LITERALES.
\usepackage[autostyle=true]{csquotes}


% ============================================================
% 3. BIBLIOGRAFÍA
% ------------------------------------------------------------
% EL ARCHIVO biblatex-<estandar>-config.tex ES RESPONSABLE DE
% CARGAR biblatex CON EL ESTILO QUE CORRESPONDA. SON NUESTROS
% ARCHIVOS, VIAJAN EN EL PAQUETE Y SE COPIAN AL PROYECTO: NO
% SE DEPENDE DE NINGÚN ESTILO REGISTRADO EN TeX Live.
%
% LA RAMA bibtex (vancouver) NO ESTÁ IMPLEMENTADA. EL XSLT
% CORTA ANTES DE GENERAR SI EL ESTILO LA REQUIERE.
% ============================================================
\input{cita-\estandar-config.tex}

% ============================================================
% Índice de autores del aparato bibliográfico
% ------------------------------------------------------------
% Lo llena biblatex: con indexing=true en el estilo, cada nombre
% citado y cada nombre de la bibliografía se indexan solos. El
% editor no escribe ningún \index.
%
% esindex ORDENA LOS APELLIDOS COMPUESTOS CASTELLANOS COMO
% CORRESPONDE: «de la Torre» SE ALFABETIZA POR Torre PERO SE
% IMPRIME CON LA PARTÍCULA. SIN ÉL, EL ORDEN POR OMISIÓN LOS
% MANDA A LA D.
%
% VA EN EL CONTRATO PORQUE ES LO QUE HACE QUE LA PIEZA
% indice_autores_bibliografia DIGA LA VERDAD. EL \makeindex ES
% DISEÑO Y VA EN LA ESTÉTICA.
% ============================================================
\usepackage{esindex}
\DeclareIndexNameFormat{default}{%
	\usebibmacro{index:name}{\esindex[names]}
	{\namepartfamily}
	{\namepartgiven}
	{\namepartprefix}
	{\namepartsuffix}}

% ============================================================
% 4. PÁGINAS Y SALTOS
% ============================================================

% ------------------------------------------------------------
% \gbclearrecto — SALTA A LA PRÓXIMA PÁGINA IMPAR
% \gbclearverso — SALTA A LA PRÓXIMA PÁGINA PAR
% ------------------------------------------------------------
% LA PÁGINA DE RELLENO LLEVA \hbox{} PORQUE UNA PÁGINA SIN
% CONTENIDO NO SE EMITE. NO SE USA UN CARÁCTER BLANCO: ESO
% DEJA TEXTO EXTRAÍBLE EN EL PDF, QUE APARECE AL COPIAR, EN
% pdftotext Y EN UN LECTOR DE PANTALLA.
% EL CONTADOR SIGUE CORRIENDO: SOLO SE OCULTA EL FOLIO.
% ------------------------------------------------------------
\makeatletter
\newcommand{\gbclearrecto}{%
  \clearpage
  \ifodd\c@page\else
    \thispagestyle{empty}\hbox{}\newpage
  \fi}
\newcommand{\gbclearverso}{%
  \clearpage
  \ifodd\c@page
    \thispagestyle{empty}\hbox{}\newpage
  \fi}
% LA BLANCA QUE GENERA \cleardoublepage SALE, POR OMISIÓN, CON
% EL FOLIO Y EL ENCABEZADO CORRIENTES. ESTE PARCHE LA DEJA
% VACÍA, QUE ES LO QUE SE ESPERA DE UNA PÁGINA EN BLANCO.
\renewcommand{\cleardoublepage}{\gbclearrecto}
\makeatother

% ------------------------------------------------------------
% gbpaginaespecial / gbpaginaimagen
% ------------------------------------------------------------
% PIEZAS SUELTAS DE PRELIMINARES: UNA DEDICATORIA, UN
% RECORDATORIO, UNA FOTO A PÁGINA ENTERA. VAN EN RECTO, SIN
% FOLIO NI ENCABEZADO, Y SU VUELTA QUEDA EN BLANCO. NO LLEVAN
% \chapter Y NO ENTRAN AL ÍNDICE.
% LA DE IMAGEN ROMPE LA CAJA DE TEXTO; LA OTRA LA RESPETA.
% ------------------------------------------------------------
\newenvironment{gbpaginaespecial}
  {\gbclearrecto\thispagestyle{empty}\begin{center}}
  {\end{center}\gbclearrecto}

\newenvironment{gbpaginaimagen}
  {\gbclearrecto\thispagestyle{empty}%
   \begingroup\centering\setlength{\parindent}{0pt}}
  {\par\endgroup\gbclearrecto}

% ------------------------------------------------------------
% gbcolofon
% ------------------------------------------------------------
% VA EN VERSO. NO PUEDE SER UN \chapter*: EN LA CLASE book,
% \chapter FUERZA EL RECTO PORQUE @openright ESTÁ ACTIVO, Y
% CUALQUIER SALTO PREVIO QUEDA ANULADO.
% ------------------------------------------------------------
\newenvironment{gbcolofon}
  {\gbclearverso\thispagestyle{empty}\justifying}
  {\clearpage}

% ============================================================
% 5. ESTRUCTURA
% ============================================================

% ------------------------------------------------------------
% \gbContadoresSinCapitulo
% ------------------------------------------------------------
% SE LLAMA UNA VEZ AL ENTRAR EN PRELIMINARES Y UNA VEZ AL
% ENTRAR EN POSLIMINARES. NO POR PIEZA.
%
% EN ESAS DOS ZONAS \chapter NO EJECUTA \refstepcounter{chapter}
% (book.cls, \@chapter, RAMA \if@mainmatter). ESE COMANDO ES LO
% ÚNICO QUE INCREMENTA EL CONTADOR DE CAPÍTULO Y, A LA VEZ,
% REINICIA LOS QUE DEPENDEN DE ÉL: section, footnote, figure,
% table Y equation. SI NO CORRE, NO PASA NINGUNA DE LAS DOS.
%
% LLEVAR chapter A CERO ACTIVA LA GUARDA \ifnum \c@chapter>\z@
% DE \thefigure, \thetable Y \theequation, QUE OMITE EL PREFIJO.
% SIN ESTO, UNA FIGURA DE POSLIMINARES SALE "Figura 1.2", CON
% EL NÚMERO DEL ÚLTIMO CAPÍTULO DEL CUERPO. VERIFICADO.
%
% LAS FIGURAS Y TABLAS NO SE REINICIAN POR PIEZA A PROPÓSITO:
% EN PRELIMINARES LA NUMERACIÓN CORRIDA (1, 2, 3) ES LA CORRECTA
% Y REINICIAR PRODUCIRÍA DOS "Figura 1" CON \label EN COLISIÓN.
% LAS NOTAS AL PIE SÍ VAN POR PIEZA: \thefootnote ES ARÁBIGO
% PELADO Y NO HAY PREFIJO QUE COLISIONE.
% ------------------------------------------------------------
\newcommand{\gbContadoresSinCapitulo}{%
  \setcounter{chapter}{0}%
  \setcounter{section}{0}%
  \setcounter{figure}{0}%
  \setcounter{table}{0}%
  \setcounter{equation}{0}%
}

% ------------------------------------------------------------
% \gbPiezaSinNumerar
% ------------------------------------------------------------
% PARA UNA PIEZA SIN NUMERAR QUE VA EN LA ZONA DE CUERPO. EL
% CASO TÍPICO SON LAS CONCLUSIONES, QUE DEBEN PRECEDER A
% \appendix PARA QUE LOS APÉNDICES NUMEREN CON LETRA.
%
% COMO ESAS PIEZAS SE EMITEN CON \chapter*, NO CORRE
% \refstepcounter{chapter} Y NO SE REINICIA NADA: SU PRIMERA
% NOTA AL PIE ARRANCA DONDE LA DEJÓ EL CAPÍTULO ANTERIOR.
% VERIFICADO.
%
% NO SE TOCA EL CONTADOR DE CAPÍTULO: EL APÉNDICE POSTERIOR LO
% NECESITA INTACTO, Y \appendix YA LO LLEVA A CERO POR SU
% CUENTA. TAMPOCO SE TOCAN figure NI table: ESAS PIEZAS NO
% LLEVAN NI CUADROS NI FIGURAS.
% ------------------------------------------------------------
% TAMBIÉN APAGA LA NUMERACIÓN DE SECCIONES. SIN ESO, LA
% PRIMERA SECCIÓN DE LAS CONCLUSIONES SALE "1.1", QUE COLISIONA
% CON LA SECCIÓN 1.1 REAL DEL CAPÍTULO 1: \thesection ES
% \thechapter.\arabic{section} Y EL CONTADOR DE CAPÍTULO SIGUE
% EN 1. VERIFICADO EN EL SUMARIO.
% LA ZONA DE APÉNDICES LA VUELVE A ENCENDER.
\newcommand{\gbPiezaSinNumerar}{%
  \setcounter{footnote}{0}%
  \setcounter{section}{0}%
  \setcounter{secnumdepth}{0}%
}

% ------------------------------------------------------------
% \gbParte{título}{subtítulo}
% ------------------------------------------------------------
% LAS PARTES NO SON ARCHIVOS: SON METADATOS DE partes_libro Y
% LAS EMITE GAMBAS EN EL main.tex. NO LLEVAN TEXTO PROPIO.
% EL SUBTÍTULO PUEDE VENIR VACÍO. AL ÍNDICE VA SOLO EL TÍTULO.
% ------------------------------------------------------------
\newcommand{\gbParte}[2]{%
  \part[#1]{#1%
    \ifstrempty{#2}{}{\\[0.4em]{\normalsize\mdseries\itshape #2}}%
  }}

% ------------------------------------------------------------
% \gbCapitulo{título}{subtítulo}
% ------------------------------------------------------------
% EQUIVALENTE PARA CAPÍTULOS CON SUBTÍTULO.
% ------------------------------------------------------------
\newcommand{\gbCapitulo}[2]{%
  \chapter[#1]{#1%
    \ifstrempty{#2}{}{\\[0.3em]{\normalsize\mdseries\itshape #2}}%
  }}

% ============================================================
% 6. BLOQUES
% ============================================================

% CITA EN BLOQUE CON ATRIBUCIÓN OPCIONAL.
% NO ES UN EPÍGRAFE: DOCBOOK LOS DISTINGUE Y EL XSLT TAMBIÉN.
\newenvironment{gbcita}
  {\begin{quote}\small}
  {\end{quote}}
\newcommand{\gbatribucion}[1]{\par\vspace{2pt}\hfill\textup{\small#1}\par}

% EPÍGRAFE. LA ATRIBUCIÓN PUEDE SER UNA CITA BIBLIOGRÁFICA:
% EL XSLT RESUELVE UN <biblioref> DE ADENTRO A \textcite.
\newcommand{\gbepigrafe}[2]{%
  \par\vspace{0.5em}%
  \begin{flushright}%
  \begin{minipage}{0.72\linewidth}%
    \small\itshape #1\par
    \vspace{2pt}\raggedleft\upshape\footnotesize #2\par
  \end{minipage}%
  \end{flushright}%
  \vspace{0.8em}\par}

% RECUADRO LATERAL. UN ARGUMENTO: info | supplementary.
\newenvironment{gbsidebar}[1]
  {\begin{tcolorbox}[colback=black!4,colframe=black!4,
     boxrule=0pt,arc=2pt,left=6pt,right=6pt,top=4pt,bottom=4pt,
     before skip=6pt,after skip=6pt]\small\sffamily}
  {\end{tcolorbox}}

% VERSO. EL XSLT YA EMITE \\ ENTRE LÍNEAS.
\newenvironment{gbverso}
  {\begin{verse}\itshape}
  {\end{verse}}

% PARLAMENTO Y LOCUTOR. DOCBOOK 5.2 BASE NO TIENE dialogue NI
% drama: EL CANÓNICO USA para role="speech" CON
% emphasis role="speaker" (RC-DB-04).
\newenvironment{gbparlamento}
  {\par\begingroup\setlength{\parindent}{0pt}\hangindent=1.2em\hangafter=1}
  {\par\endgroup}
\newcommand{\gblocutor}[1]{\textsc{\MakeLowercase{#1}}\quad}

% ============================================================
% 7. LLAMADAS DE ATENCIÓN
% ============================================================
% CADA LLAMADA DE ATENCIÓN ES UN tcolorbox CON TÍTULO FIJO.
% NO SE USA UN .style DE pgfkeys PARAMETRIZADO: EL #1 DE UN
% \newcommand DENTRO DE UN .style LO CONSUME pgfkeys Y DA
% "You already have nine parameters", QUE NO DICE NADA SOBRE
% LA CAUSA. CON UN ENTORNO POR CLASE EL PROBLEMA NO EXISTE.
\tcbset{gbadmon/.style={
  colback=black!3, colframe=black!25,
  boxrule=0.4pt, arc=1pt,
  left=6pt, right=6pt, top=4pt, bottom=4pt,
  before skip=6pt, after skip=6pt,
  fonttitle=\bfseries\sffamily\small}}

\newenvironment{gbnota}
  {\begin{tcolorbox}[gbadmon,title={Nota}]\small}{\end{tcolorbox}}
\newenvironment{gbconsejo}
  {\begin{tcolorbox}[gbadmon,title={Sugerencia}]\small}{\end{tcolorbox}}
\newenvironment{gbadvertencia}
  {\begin{tcolorbox}[gbadmon,title={Advertencia}]\small}{\end{tcolorbox}}
\newenvironment{gbimportante}
  {\begin{tcolorbox}[gbadmon,title={Importante}]\small}{\end{tcolorbox}}
\newenvironment{gbprecaucion}
  {\begin{tcolorbox}[gbadmon,title={Precaución}]\small}{\end{tcolorbox}}

% ============================================================
% 8. CÓDIGO
% ------------------------------------------------------------
% listings Y NO minted: minted EXIGE --shell-escape Y Pygments,
% Y ESO ROMPE LA AUTONOMÍA DEL PROYECTO EN OTRA MÁQUINA.
% ============================================================
\lstset{
  basicstyle=\ttfamily\footnotesize,
  breaklines=true,
  columns=fullflexible,
  frame=none,
  showstringspaces=false,
  extendedchars=true,
  inputencoding=utf8
}

% ============================================================
% 9. COLOFÓN
% ------------------------------------------------------------
% LAS FUENTES Y LA VERSIÓN SON VALORES, NO TEXTO ESCRITO. SI
% UNA EDITORIAL REEMPLAZA LA ESTÉTICA Y CAMBIA LAS FUENTES,
% EL COLOFÓN SIGUE DICIENDO LA VERDAD SIN QUE NADIE SE
% ACUERDE DE EDITARLO.
% LA FECHA ES LA DE PUBLICACIÓN Y NO LA DEL DÍA: UNA MARCA DE
% TIEMPO HARÍA QUE CADA COMPILACIÓN PRODUJERA UN DIFF AUNQUE
% NADA HUBIERA CAMBIADO.
% ============================================================
% SI LA ESTÉTICA NO LAS DECLARA, EL COLOFÓN LO DICE EN VEZ DE DAR UN
% "Undefined control sequence" QUE NO EXPLICA NADA.
\providecommand{\gbfuentetexto}{(sin declarar)}
\providecommand{\gbfuentetitulos}{(sin declarar)}
\providecommand{\gbversion}{(sin version)}

% EL RECUADRO SE DEFINE COMO ESTILO Y NO EN LÍNEA: ASÍ UNA EDITORIAL
% QUE QUIERA OTRO ASPECTO REDEFINE gbcolofoncaja SIN TOCAR EL TEXTO,
% QUE ES LO QUE EL CONTRATO GARANTIZA.
\tcbset{gbcolofoncaja/.style={
  colback=black!7,
  colframe=black!50,
  boxrule=0.7pt,
  arc=3mm,
  left=8pt, right=8pt, top=8pt, bottom=8pt,
  boxsep=0pt}}

\newcommand{\gbTextoColofon}{%
  \begin{tcolorbox}[gbcolofoncaja]
    La composición tipográfica de este libro se realizó con
    gbpublisher versión~\gbversion, bajo un modelo de fuente única que
    garantiza la trazabilidad entre el original y cada una de sus
    salidas (HTML, ePub y PDF). Las familias tipográficas utilizadas son
    \gbfuentetexto\ para el texto e \gbfuentetitulos\ para los
    títulos. El proyecto se encuentra disponible en
    \url{https://github.com/albertomoyano/gbpublisher}.
  \end{tcolorbox}}

   entornos y macros del XSLT
%   % ============================================================
% preambulo-estetica.tex
% ------------------------------------------------------------
% CAPA REEMPLAZABLE. ES LA QUE UNA EDITORIAL CAMBIA POR LA SUYA
% CUANDO QUIERE MANTENER SU IDENTIDAD VISUAL. NO DEBE DEFINIR
% NINGUNO DE LOS ENTORNOS DEL CONTRATO.
% ------------------------------------------------------------
% ESTE ARCHIVO ES EL ESQUELETO. LOS BLOQUES DE DISEÑO QUE YA
% EXISTEN EN EL preambulo.tex EN USO —titlesec, titletoc,
% estilos de página, diseño de notas al pie, newfloat,
% imakeidx, glossaries, froufrou, siunitx— SE PEGAN EN EL
% APARTADO 8. NO ESTÁN ACÁ PORQUE SON DECISIONES DE DISEÑO
% PROPIAS Y NO CORRESPONDE REESCRIBIRLAS.
% ============================================================
% LAS CUATRO PRESTACIONES
% ------------------------------------------------------------
% El derivado emite PRESTACIONES y no un «modo». Así, agregar
% un estado editorial nuevo no obliga a tocar esta capa, y cada
% bandera dice qué hace en vez de cómo se llama.
%
%   estado_libro     showframe  homeoarchy  enlaces  cruces
%   borrador             si         si         si      no
%   en_correccion        no         si         si      no
%   en_produccion        no         no         si      no
%   publicado            no         no         no      si
%
% LAS CUATRO TOCAN OPCIONES DE PAQUETE, QUE NO SE PUEDEN
% SUPERPONER DESPUES DE CARGADO EL PAQUETE. POR ESO SON
% BANDERAS DEL DERIVADO Y NO REEMPLAZOS DE TEXTO.
% LA EXCEPCION ES hyperref: SE CARGA UNA SOLA VEZ EN EL
% CONTRATO Y LOS COLORES VAN POR \hypersetup, QUE SI SE
% SUPERPONE.
% ============================================================
% CORRECCIONES APLICADAS RESPECTO DEL PREAMBULO EN USO
%   - FUERA \usepackage[utf8]{luainputenc}: LuaLaTeX CONSUME
%     UTF-8 NATIVAMENTE (SC-02).
%   - FUERA \usepackage{svg}: EXIGE Inkscape Y --shell-escape,
%     Y EN LIBROS SOLO SE USAN PNG, JPG Y PDF.
%   - \renewcommand{\footnoterule} ESTABA DEFINIDO DOS VECES
%     CON EL MISMO CUERPO. QUEDA UNO.
%   - babel LO CARGA EL CONTRATO, ANTES DE csquotes Y biblatex,
%     QUE ES EL ORDEN QUE PIDE LA DOCUMENTACION.
%   - luatexja SOLO SI EL LIBRO LO NECESITA.
% ============================================================

\usepackage{ifthen}
\usepackage{calc}

% ============================================================
% 1. GEOMETRIA
% ------------------------------------------------------------
% Los valores vienen del derivado.
%
% showframe ES OPCION DE PAQUETE Y NO SE PUEDE SUPERPONER: SE
% PASA CON \PassOptionsToPackage PARA NO DUPLICAR LA LISTA
% ENTERA DE OPCIONES EN DOS RAMAS DE UN CONDICIONAL.
%
% Hoy los cuatro margenes son iguales en los 45 formatos de la
% base, asi que inner/outer da lo mismo que centering; se usan
% igual para que el dia que se carguen asimetricos no haya que
% tocar nada. En un libro encuadernado el margen interior suele
% necesitar mas aire que el exterior.
% ============================================================
\ifgbshowframe
  \PassOptionsToPackage{showframe}{geometry}
\fi

\usepackage[
  paperwidth=\gbanchopagina,
  paperheight=\gbaltopagina,
  inner=\gbmargeninterior,
  outer=\gbmargenexterior,
  top=\gbmargensuperior,
  bottom=\gbmargeninferior,
  includehead,
  includefoot,
  headsep=14pt,
  footskip=24pt
]{geometry}

% ------------------------------------------------------------
% CRUCES DE CORTE
% ------------------------------------------------------------
% La hoja de montaje es mayor que el formato final: el derivado
% la calcula sumando el margen de cruces por lado. Las cruces
% solo necesitan lugar, y ese lugar es el mismo para cualquier
% formato.
%
% crop DECLARA LA HOJA DE MONTAJE Y geometry EL FORMATO FINAL:
% SON DOS PARES DE MEDIDAS INDEPENDIENTES, NO UNA DERIVADA DE
% LA OTRA. VERIFICADO: crop EXPANDE MACROS EN SUS OPCIONES
% keyval, ASI QUE NO HACE FALTA \crop[...] APARTE.
% ------------------------------------------------------------
\ifgbcruces
  \usepackage[width=\gbanchomontaje,height=\gbaltomontaje,cam,center]{crop}
\fi

% ============================================================
% 2. IDIOMA
% ------------------------------------------------------------
% babel lo carga el contrato, antes de csquotes y biblatex.
% Aca solo se cambian las cadenas, que si son diseno.
% ============================================================
\addto\captionsspanish{\renewcommand{\contentsname}{Sumario}}

% ============================================================
% 3. TIPOGRAFIA
% ------------------------------------------------------------
% Las fuentes son de la estetica, no de la base. Por eso cuando
% una editorial reemplaza esta capa, cambia las fuentes y el
% colofon las nombra correctamente sin que haya que tocar nada
% mas: el contrato solo las lee.
% ============================================================
\usepackage{fontspec}
\usepackage[final,babel=true]{microtype}

\def\gbfuentetexto{Libertinus Serif}
\def\gbfuentetitulos{IBM Plex Sans Condensed}
\def\gbfuentemono{IBM Plex Mono}

\setmainfont{\gbfuentetexto}[
  Numbers={OldStyle,Proportional},
  Ligatures={TeX,Common,Discretionary},Scale=1.18]
\setsansfont{\gbfuentetitulos}[
Scale=MatchLowercase,
Ligatures=TeX,
Extension=.otf,
UprightFont=*-Regular,
ItalicFont=*-Italic,
BoldFont=*-SemiBold,
BoldItalicFont=*-SemiBoldItalic]
\setmonofont{\gbfuentemono}[
Scale=0.91,
Extension=.otf,
UprightFont=*-Regular,
ItalicFont = IBMPlexMono-Italic.otf,
BoldFont = IBMPlexMono-Bold.otf,
BoldItalicFont = IBMPlexMono-BoldItalic.otf]

%\setsansfont{\gbfuentetitulos}[Scale=MatchLowercase,Ligatures=TeX]
%\setmonofont{\gbfuentemono}[Scale=0.91]

\usepackage{unicode-math}
\setmathfont{Libertinus Math}[Scale=MatchLowercase]

% SOPORTE CJK SOLO SI EL LIBRO LO NECESITA: luatexja PESA EN
% CADA COMPILACION Y CASI NINGUN LIBRO LO USA.
\ifgbcjk
  \usepackage{luatexja-fontspec}
  \setmainjfont{FandolSong}
\fi

\renewcommand{\normalsize}{\fontsize{10pt}{14pt}\selectfont}
\normalsize

% ============================================================
% 4. LISTAS, NOTAS Y FILETES
% ============================================================
% REEMPLAZA EL APARTADO 4 DE preambulo-estetica.tex
%
% QUÉ CAMBIÓ RESPECTO DEL ESQUELETO, Y POR QUÉ
% ------------------------------------------------------------
% 1. SE SACÓ hang DE footmisc
%
%    hang SANGRA EL TEXTO DE LA NOTA HACIA ADENTRO PARA HACERLE
%    LUGAR A LA MARCA. LO QUE SE BUSCA ES LO CONTRARIO: LA MARCA
%    COLGADA EN EL MARGEN Y EL TEXTO OCUPANDO EL ANCHO COMPLETO
%    DE LA CAJA. ESO LO PRODUCE \deffootnote DE scrextend.
%
%    hang Y \deffootnote GOBIERNAN LA MISMA SANGRÍA POR
%    MECANISMOS DISTINTOS, ASÍ QUE CARGAR LOS DOS DEJA EL
%    RESULTADO A MERCED DE CUÁL SE APLIQUE ÚLTIMO. QUEDA UNO.
%
%    bottom Y stable SÍ SE CONSERVAN: FIJAN LAS NOTAS AL PIE DE
%    LA CAJA Y LAS PERMITEN EN LOS TÍTULOS, Y NINGUNO TOCA LA
%    SANGRÍA.
%
% 2. SE SACARON LOS DOS \patchcmd SOBRE \@footnotetext
%
%    VERIFICADO: LOS DOS FALLAN. scrextend REDEFINE
%    \@footnotetext, ASÍ QUE CUANDO EL PARCHE BUSCA \footnotesize
%    O \@MM ESOS TOKENS YA NO ESTÁN EN LA DEFINICIÓN.
%
%    EL PRIMERO SE REEMPLAZA POR \setkomafont{footnote}, QUE ES
%    EL MECANISMO PROPIO DE KOMA Y NO DEPENDE DE PARCHEAR NADA.
%    VERIFICADO: CON \small LA MISMA NOTA OCUPA DOS PÁGINAS Y CON
%    \footnotesize UNA.
%
%    EL SEGUNDO SE REEMPLAZA POR \interfootnotelinepenalty, EL
%    PARÁMETRO DEL KERNEL QUE GOBIERNA ESE CORTE. ADVERTENCIA
%    HONESTA: NO PUDE DEMOSTRAR QUE CAMBIE NADA. UNA NOTA DE
%    CATORCE PÁRRAFOS SE PARTIÓ IGUAL CON EL PARÁMETRO Y SIN ÉL.
%    QUEDA PORQUE ES EL MECANISMO DOCUMENTADO Y NO HACE DAÑO,
%    PERO SI EL COMPORTAMIENTO CON NOTAS LARGAS NO ES EL DEL
%    PREÁMBULO LEGACY, ESTE ES EL LUGAR A REVISAR.
% ============================================================

\usepackage{enumitem}
\setlist{nosep,topsep=4pt}

% ------------------------------------------------------------
% NOTAS AL PIE
% ------------------------------------------------------------
\usepackage[bottom,stable]{footmisc}
\usepackage{scrextend}

% SEPARACIÓN ENTRE NOTAS
\setlength{\footnotesep}{10pt}

% FILETE SEPARADOR: 0,5 pt DE GROSOR Y 40% DEL ANCHO DE COLUMNA
\renewcommand{\footnoterule}{%
	\kern-3pt
	\hrule height 0.5pt width 0.4\columnwidth
	\kern6pt}

% NUMERACIÓN: ENTRE CORCHETES, PALO SECO, CUERPO scriptsize
\renewcommand*{\thefootnote}{\scriptsize\sffamily{[\arabic{footnote}]}}

% CUERPO DE LA NOTA
\setkomafont{footnote}{\small}

% ------------------------------------------------------------
% MARCA COLGADA EN EL MARGEN
% ------------------------------------------------------------
% \llap SACA LA MARCA FUERA DE LA CAJA HACIA LA IZQUIERDA, ASÍ EL
% TEXTO DE LA NOTA EMPIEZA EN EL MARGEN Y OCUPA EL ANCHO COMPLETO.
%
% LOS TRES ARGUMENTOS DE \deffootnote SON: SANGRÍA DE LA PRIMERA
% LÍNEA, SANGRÍA DE LAS SIGUIENTES, Y CÓMO SE COMPONE LA MARCA.
% ------------------------------------------------------------
\newcommand*\gbmarcanotaespacio{0em}

\deffootnote{\gbmarcanotaespacio}
{\parindent}
{%
	\makebox[\gbmarcanotaespacio][r]{\llap{\thefootnotemark\quad}}%
}

% ------------------------------------------------------------
% NOTAS LARGAS PARTIDAS ENTRE PÁGINAS
% ------------------------------------------------------------
% Penalización de partir una nota entre dos páginas. Por omisión
% LaTeX pone 100 en la clase book, y el valor alto que usa el
% kernel en otros contextos hace que una nota muy larga desplace
% texto en lugar de partirse.
%
% VER LA ADVERTENCIA DEL ENCABEZADO: NO PUDE MEDIR DIFERENCIA.
% ------------------------------------------------------------
\interfootnotelinepenalty=100

% ------------------------------------------------------------
% RESTO DEL APARTADO
% ------------------------------------------------------------
\usepackage[labelfont=bf,font=small,labelsep=period,format=plain]{caption}
\usepackage{needspace}

\setcounter{tocdepth}{1}
\setcounter{secnumdepth}{3}
\renewcommand{\thepart}{\arabic{part}}

\raggedbottom
\clubpenalty=10000
\widowpenalty=10000

% ============================================================
% 5. MARCAS TIPOGRAFICAS DE TRABAJO
% ------------------------------------------------------------
% impnattypo marca homeoarquias y problemas de principio y fin
% de linea. La opcion draft es opcion de paquete y no se puede
% superponer: por eso la carga va condicionada y no comentada
% a mano.
% ============================================================
\ifgbhomeoarchy
  \IfFileExists{impnattypo.sty}{%
    \usepackage[hyphenation,homeoarchy,homeoarchywordcolor=orange,
                homeoarchycharcolor=orange,draft]{impnattypo}}{%
    \PackageWarning{gbpublisher}{impnattypo no instalado:
      se pierden las marcas tipograficas de trabajo}}
  \overfullrule=5pt
\fi

% ============================================================
% 6. COLOR DE LOS ENLACES
% ------------------------------------------------------------
% Solo el color. hyperref lo carga preambulo-hyperref.tex, que
% va después de esta capa y lee \ifgbenlaces para decidir entre
% enlaces coloreados y hidelinks.
% ============================================================
\usepackage{xcolor}
\definecolor{gbenlace}{RGB}{160,0,120}

% ============================================================
% 7. CONSTANCIA EN EL LOG
% ------------------------------------------------------------
% Deja escrito en el .log con que prestaciones se compuso.
% Cuando un PDF sale distinto de lo esperado, esta linea dice
% por que sin tener que abrir el derivado.
% ============================================================
\typeout{[gbpublisher] prestaciones:
  showframe=\ifgbshowframe si\else no\fi,
  homeoarchy=\ifgbhomeoarchy si\else no\fi,
  enlaces=\ifgbenlaces si\else no\fi,
  cruces=\ifgbcruces si\else no\fi}

% ============================================================
% 8. AQUI VAN LOS BLOQUES DE DISENO DEL PREAMBULO EN USO
% ------------------------------------------------------------
%   - titlesec y los estilos de titulo
%   - titletoc: partes, capitulos y secciones del sumario
%   - titleps: estilos de pagina y folios corridos
%     (\gbtituloabreviado ES LO QUE CORRESPONDE EN EL FOLIO
%      CORRIDO, Y HOY EL PREAMBULO NO LO CONSUME)
%   - newfloat: el entorno flotante "imagen"
%   - imakeidx + xindy: names (lo llena esindex desde biblatex),
%     onomastico y concepto (los carga el editor)
%   - glossaries: siglas y glosario
%   - siunitx, froufrou, qrcode, pdfpages, rotating
% ============================================================

% Cita con cambio de márgenes y tamaño tipografico
\renewenvironment{quote}
{\normalsize\list{}{\sf\leftmargin=14pt \rightmargin=0pt}%
	\item\relax}
{\endlist}

% Configuración de los índices
\usepackage[xindy]{imakeidx}
\makeindex[name=names, title={}, intoc=false]
\makeindex[name=onomastico, title={}, intoc=false]
\makeindex[name=concepto, title={}, intoc=false]


% ============================================================
% PERSONALIZACIÓN DE LOS ÍNDICES DE PARTES, CAPÍTULOS Y SECCIONES
% ------------------------------------------------------------
% Paquete requerido:
\usepackage{titletoc}
%
% OBJETIVO GENERAL:
% - Controlar la apariencia de los índices (Tabla de Contenidos)
%   para partes, capítulos y secciones.
% - Ajustar sangrías, tipografía, alineación y estilo de los
%   números de página.
% ============================================================

% ------------------------------------------------------------
% PARTES
% ------------------------------------------------------------
% \titlecontents{part}[<left>]{<above-code>}{<numbered-entry-format>}
%                        {<numberless-entry-format>}{<filler-page>}
%
% Argumentos usados:
%   [0em]                    -> Sangría izquierda de la entrada.
%   {\addvspace{5pt}\sf\bfseries\normalsize\selectfont\filright}
%                             -> Espacio vertical antes de la entrada (5pt),
%                                fuente sans serif, negrita, tamaño normal,
%                                alineación a la derecha (\filright).
%   {\contentslabel[\thecontentslabel]{2.5pc}}
%                             -> Formato de la etiqueta numérica (ej. Parte I)
%                                con ancho de 2.5pc.
%   {}                        -> Formato de entrada sin número.
%   {}                        -> No hay reglas ni página aquí.
\titlecontents{part}[0em]
{\addvspace{5pt}\sf\bfseries\normalsize\selectfont\filright}
{\contentslabel[\thecontentslabel]{2.5pc}}
{}
{}

% ------------------------------------------------------------
% CAPÍTULOS
% ------------------------------------------------------------
% Formato general:
%   - Sangría: 1.5pc
%   - Espacio vertical: 0.4em antes de cada capítulo
%   - Fuente: sans serif, alineación a la derecha (\filright)
%   - Número del capítulo alineado con margen de 1.5pc
%   - Separador de puntos hasta el número de página (\titlerule)
%
% Argumentos:
%   [1.5pc] -> sangría izquierda del capítulo
%   {\addvspace{.4em}\sf\selectfont\filright} -> formato general de la entrada
%   {\contentslabel{1.5pc}} -> ancho reservado para número
%   {\hspace*{-1.5pc}} -> corrección si no tiene número
%   {\titlerule*[1pc]{.}\contentspage[\hspace*{-4pc} {\rm\small\thecontentspage}]}
%      -> regla de puntos de relleno hasta el número de página,
%         página en tamaño pequeño y ligeramente desplazada
\titlecontents{chapter}[1.5pc]
{\addvspace{.4em}\sf\selectfont\filright}
{\contentslabel{1.5pc}}
{\hspace*{-1.5pc}}%
{\titlerule*[1pc]{.}\contentspage[\hspace*{-4pc} {\rm\small\thecontentspage}]}%
[]

% ------------------------------------------------------------
% SECCIONES
% ------------------------------------------------------------
% Formato general:
%   - Sangría: 4.5pc
%   - Fuente: \small, alineación a la derecha
%   - Número de sección: 2.5pc
%   - Ajuste de sangría si no hay número: -2.5pc
%   - Separador de puntos hasta el número de página (\titlerule)
\titlecontents{section}[4.5pc]
{\small\filright}                  % Formato general
{\contentslabel{2.5pc}}            % Número con ancho 2.5pc
{\hspace*{-2.5pc}}                 % Si no tiene número
{\titlerule*[1pc]{.}\contentspage} % Regla de puntos hasta la página


% ============================================================
% CONFIGURACIÓN DE ESTILOS DE TÍTULOS CON titlesec
% ------------------------------------------------------------
% Paquete requerido:
\usepackage[sf,bf,compact,topmarks,calcwidth,pagestyles,clearempty,newlinetospace]{titlesec}
%
% OBJETIVO GENERAL:
% - Personalizar la apariencia de partes, capítulos, secciones,
%   subsecciones, subsubsecciones, párrafos y subpárrafos.
% - Controlar fuente, tamaño, alineación, sangrías y espaciado.
% ============================================================

% ============================================================
% DISEÑO DE PARTES
% ------------------------------------------------------------
% Redefinición de \@part y \@spart para controlar:
% - Numeración de partes (si corresponde)
% - Entrada en el TOC
% - Centrado del título y subtítulo de la parte
% - Tamaño de fuente y estilo tipográfico
%
% \@part[#1]#2:
%   - #1: título corto para TOC y encabezado
%   - #2: título completo para mostrar en la página
%   - Centrado, fuente sans serif (\sf), \LARGE para nombre de parte
%     y \Large para título completo
%
% \@spart#1:
%   - Para partes no numeradas
%   - Solo muestra el título centrado, \LARGE, sans serif
%
\makeatletter
\def\@part[#1]#2{%
	\ifnum \c@secnumdepth >-2\relax
	\refstepcounter{part}%
	\addcontentsline{toc}{part}{Parte \thepart\hspace{1em}#1}%
	\else
	\addcontentsline{toc}{part}{#1}%
	\fi
	\markboth{}{}%
	{\centering
		\interlinepenalty \@M
		\ifnum \c@secnumdepth >-2\relax
		\sf\LARGE\selectfont \partname~\thepart
		\par\vskip 20\p@%
		\fi
		\sf\Large\selectfont
		#2\par}%
	\@endpart
}

\def\@spart#1{%
	\markboth{}{}%
	{\centering
		\interlinepenalty \@M
		\sf\LARGE\selectfont
		#1\par}%
	\@endpart
}
\makeatother

% ============================================================
% DISEÑO DE CAPÍTULO
% ------------------------------------------------------------
% \titleformat{\chapter}[display]{...}{...}{sep}{before-code}[after-code]
% - [display]: título en bloque separado (no inline)
% - \Large: tamaño base del título
% - \vspace{-2.5cm}: ajusta altura vertical (sube el título)
% - \centering: centrado
% - \textsc{\MakeLowercase\chaptertitlename}: 'Capítulo' en versalitas
% - \thechapter: número del capítulo
% - 1.5cm: separación entre número y texto
% - \filright\sf\LARGE: alineación a la derecha, fuente sans serif, tamaño LARGER
\titleformat{\chapter}[display]
{\Large}
{\centering{\textsc{\MakeLowercase\chaptertitlename}}~\thechapter}
{1.5cm}
{\filright\sf\LARGE}
[]
% ------------------------------------------------------------
% ESPACIADO DEL TÍTULO DE CAPÍTULO
% ------------------------------------------------------------
% \titlespacing*{\chapter}{izquierda}{antes}{después}
%
% ANTES = 0pt HACE QUE «capítulo N» ARRANQUE EN EL BORDE SUPERIOR
% DE LA CAJA. LA CLASE book RESERVA 50pt POR OMISIÓN.
%
% LA VERSIÓN CON ASTERISCO SUPRIME LA SANGRÍA DEL PRIMER PÁRRAFO
% DESPUÉS DEL TÍTULO.
%
% NO SE USA \vspace NEGATIVO DENTRO DEL \titleformat: ESO EMPUJA
% HACIA ARRIBA DESPUÉS DE QUE LaTeX YA RESERVÓ EL ESPACIO, Y EL
% RESULTADO DEPENDE DEL CONTEXTO. titlesec TIENE EL PARÁMETRO
% PARA DECIRLO DIRECTO.
% ------------------------------------------------------------
\titlespacing*{\chapter}{0pt}{0pt}{30pt}

% ============================================================
% DISEÑO DE SECCIÓN
% ------------------------------------------------------------
% Fuente: sans serif, negrita, tamaño grande (\large), alineación a la derecha
\titleformat{\section}[hang]
{\sf\bfseries\large\raggedright}
{\thesection}{.5em}{}[]
\titlespacing{\section}
{\parindent}{18pt plus 0.5pt minus 0.5pt}{6.75pt}

% ============================================================
% DISEÑO DE SUBSECCIÓN
% ------------------------------------------------------------
\titleformat{\subsection}[hang]
{\sf\large\raggedright}
{\thesubsection}{.5em}{}[]
\titlespacing{\subsection}
{\parindent}{18pt plus 0.5pt minus 0.5pt}{6.75pt}

% ============================================================
% DISEÑO DE SUBSUBSECCIÓN
% ------------------------------------------------------------
\titleformat{\subsubsection}[hang]
{\rm\bfseries\normalsize\raggedright}
{\thesubsubsection}{.5em}{}[]
\titlespacing{\subsubsection}
{\parindent}{18pt plus 0.5pt minus 0.5pt}{6.75pt}

% ============================================================
% DISEÑO DE PÁRRAFO
% ------------------------------------------------------------
% \paragraph: run-in (alineado inline, no bloque)
% - Negrita
% - Sin número
% - Raya al final del título (---)
\titlespacing{\paragraph}
{0pt}                % margen izquierda
{6.75pt plus 0.5pt minus 0.5pt} % espacio antes
{0pt}                % espacio después
\titleformat{\paragraph}[runin]
{\bfseries}{}{0em}{}[\mbox{ --- }]

% ============================================================
% DISEÑO DE SUBPÁRRAFO
% ------------------------------------------------------------
% \subparagraph: estilo colapsado, centrado
% - Fuente sans serif, negrita, tamaño \large
\titlespacing{\subparagraph}
{\parindent}{18pt plus 0.5pt minus 0.5pt}{6.75pt}
\titleformat{\subparagraph}[hang]
{\sf\bfseries\large\centering}
{\thesubparagraph}{.5em}{}[]

   diseño, reemplazable
%   % ============================================================
% preambulo-hyperref.tex
% ------------------------------------------------------------
% CUARTA Y ÚLTIMA CAPA DEL PREÁMBULO. VA DESPUÉS DE LA ESTÉTICA.
%
%   % ============================================================
% preambulo-proyecto.tex
% ------------------------------------------------------------
% ARCHIVO DERIVADO — GENERADO POR gbpublisher
% NO EDITAR A MANO: SE REESCRIBE EN CADA GENERACIÓN DEL
% main.tex. LOS CAMBIOS SE HACEN EN LA BASE.
%
% SÍ SE VERSIONA: EL COMMIT ES LA GARANTÍA DE TRAZA DE QUÉ SE
% PUBLICÓ. POR ESO NO LLEVA MARCA DE TIEMPO.
% ------------------------------------------------------------
% ESTA CAPA SOLO DEFINE VALORES. NO CARGA NINGÚN PAQUETE.
% VA PRIMERA PORQUE geometry Y fontspec NECESITAN LOS VALORES
% EN EL MOMENTO DEL \usepackage.
% ============================================================

% --- PRESTACIONES DE SALIDA (derivadas de estado_libro) ---
\newif\ifgbshowframe    \gbshowframefalse
\newif\ifgbhomeoarchy   \gbhomeoarchytrue
\newif\ifgbenlaces      \gbenlacestrue
\newif\ifgbcruces       \gbcrucesfalse

% --- FORMATO DE PÁGINA (formatos_pdf id=32) ---
\def\gbanchopagina{15.50cm}
\def\gbaltopagina{23.00cm}
\def\gbmargeninterior{2.00cm}
\def\gbmargenexterior{2.00cm}
\def\gbmargensuperior{2.00cm}
\def\gbmargeninferior{2.00cm}
\def\gbnombreformato{Estudio 2A / 04}
\def\gbanchomontaje{18.00truecm}
\def\gbaltomontaje{25.50truecm}

% --- IDIOMA ---
\def\gbidiomas{main=spanish}
\def\gbidiomaprincipal{spanish}
\newif\ifgbcjk   \gbcjkfalse

% --- BIBLIOGRAFÍA ---
\def\estandar{apa}
\newif\ifgbbibporcapitulo   \gbbibporcapitulotrue
\def\gbarchivobib{./referencias/ref-l-DUEK.bib}

% --- METADATOS DEL PDF ---
\def\gbtitulo{Un tiempo diferente}
\def\gbsubtitulo{Vida cotidiana y pandemia}
\def\gbtituloabreviado{}
\def\gbautores{Moguillansky, Marina (comp.); Duek, Carolina (comp.)}
\def\gbeditorial{Ediciones Imago Mundi}
\def\gbciudad{Buenos Aires}
\def\gbanio{2026}
\def\gbisbn{978-950-793-487-2}
\def\gbisbnelectronico{}
\def\gbdoi{}
\def\gblicencia{}
\def\gburllicencia{}
\def\gbpalabrasclave{}

% --- VERSIÓN DE gbpublisher ---
\def\gbversion{0.0.332}
   valores de la base
%   % ============================================================
% preambulo-contrato.tex
% ------------------------------------------------------------
% CONTRATO VERSIÓN 1
% ------------------------------------------------------------
% DEFINE TODO LO QUE docbook-to-latex.xsl PUEDE EMITIR EN LOS
% FRAGMENTOS DE CAPÍTULO. ES LA CAPA QUE UNA EDITORIAL
% CONSERVA CUANDO REEMPLAZA LA ESTÉTICA: SI FALTA UNA DE
% ESTAS DEFINICIONES, EL LIBRO NO COMPILA; SI ESTÁ DEFINIDA
% DISTINTO, COMPILA Y SALE MAL.
%
% CUANDO SE AGREGUE VOCABULARIO NUEVO, SUBIR LA VERSIÓN DEL
% CONTRATO. EL main.tex DECLARA LA QUE NECESITA Y ESTA CAPA
% LA VERIFICA: ASÍ UN PREÁMBULO VIEJO FRENA EN LA PRIMERA
% LÍNEA CON UN MENSAJE ÚTIL, EN VEZ DE DAR UN
% "Undefined control sequence" EN LA PÁGINA 140.
% ------------------------------------------------------------
% DEPENDE DE preambulo-proyecto.tex, QUE VA ANTES.
% ============================================================

\def\gbContratoVersion{1}

% ============================================================
% \gbRequiereContrato{N}
% PROPÓSITO : EL main.tex DECLARA QUÉ VERSIÓN NECESITA.
% PARÁMETRO : N — versión requerida
% ============================================================
\newcommand{\gbRequiereContrato}[1]{%
  \ifnum#1>\gbContratoVersion\relax
    \PackageError{gbpublisher}{%
      Este preambulo implementa el contrato \gbContratoVersion\space
      y el libro requiere el #1}{%
      Actualizar preambulo-contrato.tex.}%
  \fi}

% ============================================================
% HOJAS DE CORTESÍA
% ------------------------------------------------------------
% Dos páginas en blanco al inicio —una hoja— que todo libro
% lleva. Entran en la paginación: son i y ii, y el sumario
% empieza en iii.
%
% \hbox{} ES IMPRESCINDIBLE: UNA PÁGINA SIN NINGÚN MATERIAL NO
% SE EMITE, ASÍ QUE UN \newpage SOBRE UNA PÁGINA VACÍA NO
% PRODUCE NADA Y EL CONTADOR NO AVANZA. LA CAJA VACÍA DA EL
% MÍNIMO DE CONTENIDO QUE OBLIGA A COMPONERLA.
%
% \thispagestyle{empty} LES QUITA EL FOLIO: SE CUENTAN PERO NO
% SE IMPRIME EL NÚMERO.
% ============================================================
\newcommand{\gbHojasCortesia}{%
	\thispagestyle{empty}\hbox{}\newpage
	\thispagestyle{empty}\hbox{}\newpage}

% ============================================================
% ENCABEZADO DE PIEZA GENERADA
% ------------------------------------------------------------
% Para las piezas sin cuerpo que llevan título propio. El título
% sale de capitulos.titulo_capitulo, que es donde el editor lo
% escribe, y no del preámbulo.
%
% \markboth actualiza el folio corrido, que \chapter* no toca.
% ============================================================
\newcommand{\gbEncabezadoPieza}[1]{%
	\chapter*{#1}%
	\addcontentsline{toc}{chapter}{#1}%
	\markboth{\MakeUppercase{#1}}{\MakeUppercase{#1}}}

% ============================================================
% ÍNDICE CON SU ENCABEZADO
% ------------------------------------------------------------
% \printindex HACE UN \clearpage INTERNO AUNQUE EL ÍNDICE SE
% DECLARE CON title={}. SIN NEUTRALIZARLO, EL ENCABEZADO QUEDA
% SOLO EN UNA PÁGINA Y LA LISTA EMPIEZA EN LA SIGUIENTE.
%
% EL \let VA DENTRO DE UN GRUPO: FUERA DE ÉL \clearpage TIENE QUE
% SEGUIR FUNCIONANDO, Y EL COLOFÓN QUE VIENE DESPUÉS LO NECESITA.
% ============================================================
\newcommand{\gbIndice}[2]{%
	\gbEncabezadoPieza{#1}%
	\begingroup
	\let\clearpage\relax
	\raggedright
	\small
	\printindex[#2]%
	\endgroup}

% ============================================================
% 1. PAQUETES QUE EL CONTRATO NECESITA
% ------------------------------------------------------------
% SOLO LO QUE HACE FALTA PARA QUE UN FRAGMENTO COMPILE. LA
% ESTÉTICA CARGA LO SUYO APARTE. TODOS ESTÁN EN TeX Live Y
% NINGUNO PIDE --shell-escape.
% ============================================================
\usepackage{etoolbox}
\usepackage{graphicx}
\usepackage{booktabs}
\usepackage{array}
\usepackage{listings}
\usepackage{ragged2e}
\usepackage{xcolor}
\usepackage[most]{tcolorbox}

% ------------------------------------------------------------
% IDIOMA
% ------------------------------------------------------------
% babel VA EN EL CONTRATO Y NO EN LA ESTÉTICA, POR DOS RAZONES.
% LA PRIMERA ES DE ORDEN: biblatex PIDE CARGARSE DESPUÉS DE
% babel Y csquotes, Y LA ESTÉTICA VA DESPUÉS DE ESTA CAPA.
% LA SEGUNDA ES DE FONDO: EL IDIOMA DECIDE LAS COMILLAS QUE
% PRODUCE \enquote, LA PARTICIÓN DE PALABRAS Y LAS CADENAS DE
% LA BIBLIOGRAFÍA. SON GARANTÍAS DEL CONTRATO, NO DECISIONES
% DE DISEÑO: UNA EDITORIAL QUE REEMPLAZA LA ESTÉTICA NO DEBE
% PODER CAMBIAR EL IDIOMA DEL LIBRO.
\usepackage[\gbidiomas]{babel}
\frenchspacing

% COMILLAS SEGÚN EL IDIOMA ACTIVO: EN ESPAÑOL DA « ».
% EL XSLT EMITE \enquote Y NUNCA COMILLAS LITERALES.
\usepackage[autostyle=true]{csquotes}


% ============================================================
% 3. BIBLIOGRAFÍA
% ------------------------------------------------------------
% EL ARCHIVO biblatex-<estandar>-config.tex ES RESPONSABLE DE
% CARGAR biblatex CON EL ESTILO QUE CORRESPONDA. SON NUESTROS
% ARCHIVOS, VIAJAN EN EL PAQUETE Y SE COPIAN AL PROYECTO: NO
% SE DEPENDE DE NINGÚN ESTILO REGISTRADO EN TeX Live.
%
% LA RAMA bibtex (vancouver) NO ESTÁ IMPLEMENTADA. EL XSLT
% CORTA ANTES DE GENERAR SI EL ESTILO LA REQUIERE.
% ============================================================
\input{cita-\estandar-config.tex}

% ============================================================
% Índice de autores del aparato bibliográfico
% ------------------------------------------------------------
% Lo llena biblatex: con indexing=true en el estilo, cada nombre
% citado y cada nombre de la bibliografía se indexan solos. El
% editor no escribe ningún \index.
%
% esindex ORDENA LOS APELLIDOS COMPUESTOS CASTELLANOS COMO
% CORRESPONDE: «de la Torre» SE ALFABETIZA POR Torre PERO SE
% IMPRIME CON LA PARTÍCULA. SIN ÉL, EL ORDEN POR OMISIÓN LOS
% MANDA A LA D.
%
% VA EN EL CONTRATO PORQUE ES LO QUE HACE QUE LA PIEZA
% indice_autores_bibliografia DIGA LA VERDAD. EL \makeindex ES
% DISEÑO Y VA EN LA ESTÉTICA.
% ============================================================
\usepackage{esindex}
\DeclareIndexNameFormat{default}{%
	\usebibmacro{index:name}{\esindex[names]}
	{\namepartfamily}
	{\namepartgiven}
	{\namepartprefix}
	{\namepartsuffix}}

% ============================================================
% 4. PÁGINAS Y SALTOS
% ============================================================

% ------------------------------------------------------------
% \gbclearrecto — SALTA A LA PRÓXIMA PÁGINA IMPAR
% \gbclearverso — SALTA A LA PRÓXIMA PÁGINA PAR
% ------------------------------------------------------------
% LA PÁGINA DE RELLENO LLEVA \hbox{} PORQUE UNA PÁGINA SIN
% CONTENIDO NO SE EMITE. NO SE USA UN CARÁCTER BLANCO: ESO
% DEJA TEXTO EXTRAÍBLE EN EL PDF, QUE APARECE AL COPIAR, EN
% pdftotext Y EN UN LECTOR DE PANTALLA.
% EL CONTADOR SIGUE CORRIENDO: SOLO SE OCULTA EL FOLIO.
% ------------------------------------------------------------
\makeatletter
\newcommand{\gbclearrecto}{%
  \clearpage
  \ifodd\c@page\else
    \thispagestyle{empty}\hbox{}\newpage
  \fi}
\newcommand{\gbclearverso}{%
  \clearpage
  \ifodd\c@page
    \thispagestyle{empty}\hbox{}\newpage
  \fi}
% LA BLANCA QUE GENERA \cleardoublepage SALE, POR OMISIÓN, CON
% EL FOLIO Y EL ENCABEZADO CORRIENTES. ESTE PARCHE LA DEJA
% VACÍA, QUE ES LO QUE SE ESPERA DE UNA PÁGINA EN BLANCO.
\renewcommand{\cleardoublepage}{\gbclearrecto}
\makeatother

% ------------------------------------------------------------
% gbpaginaespecial / gbpaginaimagen
% ------------------------------------------------------------
% PIEZAS SUELTAS DE PRELIMINARES: UNA DEDICATORIA, UN
% RECORDATORIO, UNA FOTO A PÁGINA ENTERA. VAN EN RECTO, SIN
% FOLIO NI ENCABEZADO, Y SU VUELTA QUEDA EN BLANCO. NO LLEVAN
% \chapter Y NO ENTRAN AL ÍNDICE.
% LA DE IMAGEN ROMPE LA CAJA DE TEXTO; LA OTRA LA RESPETA.
% ------------------------------------------------------------
\newenvironment{gbpaginaespecial}
  {\gbclearrecto\thispagestyle{empty}\begin{center}}
  {\end{center}\gbclearrecto}

\newenvironment{gbpaginaimagen}
  {\gbclearrecto\thispagestyle{empty}%
   \begingroup\centering\setlength{\parindent}{0pt}}
  {\par\endgroup\gbclearrecto}

% ------------------------------------------------------------
% gbcolofon
% ------------------------------------------------------------
% VA EN VERSO. NO PUEDE SER UN \chapter*: EN LA CLASE book,
% \chapter FUERZA EL RECTO PORQUE @openright ESTÁ ACTIVO, Y
% CUALQUIER SALTO PREVIO QUEDA ANULADO.
% ------------------------------------------------------------
\newenvironment{gbcolofon}
  {\gbclearverso\thispagestyle{empty}\justifying}
  {\clearpage}

% ============================================================
% 5. ESTRUCTURA
% ============================================================

% ------------------------------------------------------------
% \gbContadoresSinCapitulo
% ------------------------------------------------------------
% SE LLAMA UNA VEZ AL ENTRAR EN PRELIMINARES Y UNA VEZ AL
% ENTRAR EN POSLIMINARES. NO POR PIEZA.
%
% EN ESAS DOS ZONAS \chapter NO EJECUTA \refstepcounter{chapter}
% (book.cls, \@chapter, RAMA \if@mainmatter). ESE COMANDO ES LO
% ÚNICO QUE INCREMENTA EL CONTADOR DE CAPÍTULO Y, A LA VEZ,
% REINICIA LOS QUE DEPENDEN DE ÉL: section, footnote, figure,
% table Y equation. SI NO CORRE, NO PASA NINGUNA DE LAS DOS.
%
% LLEVAR chapter A CERO ACTIVA LA GUARDA \ifnum \c@chapter>\z@
% DE \thefigure, \thetable Y \theequation, QUE OMITE EL PREFIJO.
% SIN ESTO, UNA FIGURA DE POSLIMINARES SALE "Figura 1.2", CON
% EL NÚMERO DEL ÚLTIMO CAPÍTULO DEL CUERPO. VERIFICADO.
%
% LAS FIGURAS Y TABLAS NO SE REINICIAN POR PIEZA A PROPÓSITO:
% EN PRELIMINARES LA NUMERACIÓN CORRIDA (1, 2, 3) ES LA CORRECTA
% Y REINICIAR PRODUCIRÍA DOS "Figura 1" CON \label EN COLISIÓN.
% LAS NOTAS AL PIE SÍ VAN POR PIEZA: \thefootnote ES ARÁBIGO
% PELADO Y NO HAY PREFIJO QUE COLISIONE.
% ------------------------------------------------------------
\newcommand{\gbContadoresSinCapitulo}{%
  \setcounter{chapter}{0}%
  \setcounter{section}{0}%
  \setcounter{figure}{0}%
  \setcounter{table}{0}%
  \setcounter{equation}{0}%
}

% ------------------------------------------------------------
% \gbPiezaSinNumerar
% ------------------------------------------------------------
% PARA UNA PIEZA SIN NUMERAR QUE VA EN LA ZONA DE CUERPO. EL
% CASO TÍPICO SON LAS CONCLUSIONES, QUE DEBEN PRECEDER A
% \appendix PARA QUE LOS APÉNDICES NUMEREN CON LETRA.
%
% COMO ESAS PIEZAS SE EMITEN CON \chapter*, NO CORRE
% \refstepcounter{chapter} Y NO SE REINICIA NADA: SU PRIMERA
% NOTA AL PIE ARRANCA DONDE LA DEJÓ EL CAPÍTULO ANTERIOR.
% VERIFICADO.
%
% NO SE TOCA EL CONTADOR DE CAPÍTULO: EL APÉNDICE POSTERIOR LO
% NECESITA INTACTO, Y \appendix YA LO LLEVA A CERO POR SU
% CUENTA. TAMPOCO SE TOCAN figure NI table: ESAS PIEZAS NO
% LLEVAN NI CUADROS NI FIGURAS.
% ------------------------------------------------------------
% TAMBIÉN APAGA LA NUMERACIÓN DE SECCIONES. SIN ESO, LA
% PRIMERA SECCIÓN DE LAS CONCLUSIONES SALE "1.1", QUE COLISIONA
% CON LA SECCIÓN 1.1 REAL DEL CAPÍTULO 1: \thesection ES
% \thechapter.\arabic{section} Y EL CONTADOR DE CAPÍTULO SIGUE
% EN 1. VERIFICADO EN EL SUMARIO.
% LA ZONA DE APÉNDICES LA VUELVE A ENCENDER.
\newcommand{\gbPiezaSinNumerar}{%
  \setcounter{footnote}{0}%
  \setcounter{section}{0}%
  \setcounter{secnumdepth}{0}%
}

% ------------------------------------------------------------
% \gbParte{título}{subtítulo}
% ------------------------------------------------------------
% LAS PARTES NO SON ARCHIVOS: SON METADATOS DE partes_libro Y
% LAS EMITE GAMBAS EN EL main.tex. NO LLEVAN TEXTO PROPIO.
% EL SUBTÍTULO PUEDE VENIR VACÍO. AL ÍNDICE VA SOLO EL TÍTULO.
% ------------------------------------------------------------
\newcommand{\gbParte}[2]{%
  \part[#1]{#1%
    \ifstrempty{#2}{}{\\[0.4em]{\normalsize\mdseries\itshape #2}}%
  }}

% ------------------------------------------------------------
% \gbCapitulo{título}{subtítulo}
% ------------------------------------------------------------
% EQUIVALENTE PARA CAPÍTULOS CON SUBTÍTULO.
% ------------------------------------------------------------
\newcommand{\gbCapitulo}[2]{%
  \chapter[#1]{#1%
    \ifstrempty{#2}{}{\\[0.3em]{\normalsize\mdseries\itshape #2}}%
  }}

% ============================================================
% 6. BLOQUES
% ============================================================

% CITA EN BLOQUE CON ATRIBUCIÓN OPCIONAL.
% NO ES UN EPÍGRAFE: DOCBOOK LOS DISTINGUE Y EL XSLT TAMBIÉN.
\newenvironment{gbcita}
  {\begin{quote}\small}
  {\end{quote}}
\newcommand{\gbatribucion}[1]{\par\vspace{2pt}\hfill\textup{\small#1}\par}

% EPÍGRAFE. LA ATRIBUCIÓN PUEDE SER UNA CITA BIBLIOGRÁFICA:
% EL XSLT RESUELVE UN <biblioref> DE ADENTRO A \textcite.
\newcommand{\gbepigrafe}[2]{%
  \par\vspace{0.5em}%
  \begin{flushright}%
  \begin{minipage}{0.72\linewidth}%
    \small\itshape #1\par
    \vspace{2pt}\raggedleft\upshape\footnotesize #2\par
  \end{minipage}%
  \end{flushright}%
  \vspace{0.8em}\par}

% RECUADRO LATERAL. UN ARGUMENTO: info | supplementary.
\newenvironment{gbsidebar}[1]
  {\begin{tcolorbox}[colback=black!4,colframe=black!4,
     boxrule=0pt,arc=2pt,left=6pt,right=6pt,top=4pt,bottom=4pt,
     before skip=6pt,after skip=6pt]\small\sffamily}
  {\end{tcolorbox}}

% VERSO. EL XSLT YA EMITE \\ ENTRE LÍNEAS.
\newenvironment{gbverso}
  {\begin{verse}\itshape}
  {\end{verse}}

% PARLAMENTO Y LOCUTOR. DOCBOOK 5.2 BASE NO TIENE dialogue NI
% drama: EL CANÓNICO USA para role="speech" CON
% emphasis role="speaker" (RC-DB-04).
\newenvironment{gbparlamento}
  {\par\begingroup\setlength{\parindent}{0pt}\hangindent=1.2em\hangafter=1}
  {\par\endgroup}
\newcommand{\gblocutor}[1]{\textsc{\MakeLowercase{#1}}\quad}

% ============================================================
% 7. LLAMADAS DE ATENCIÓN
% ============================================================
% CADA LLAMADA DE ATENCIÓN ES UN tcolorbox CON TÍTULO FIJO.
% NO SE USA UN .style DE pgfkeys PARAMETRIZADO: EL #1 DE UN
% \newcommand DENTRO DE UN .style LO CONSUME pgfkeys Y DA
% "You already have nine parameters", QUE NO DICE NADA SOBRE
% LA CAUSA. CON UN ENTORNO POR CLASE EL PROBLEMA NO EXISTE.
\tcbset{gbadmon/.style={
  colback=black!3, colframe=black!25,
  boxrule=0.4pt, arc=1pt,
  left=6pt, right=6pt, top=4pt, bottom=4pt,
  before skip=6pt, after skip=6pt,
  fonttitle=\bfseries\sffamily\small}}

\newenvironment{gbnota}
  {\begin{tcolorbox}[gbadmon,title={Nota}]\small}{\end{tcolorbox}}
\newenvironment{gbconsejo}
  {\begin{tcolorbox}[gbadmon,title={Sugerencia}]\small}{\end{tcolorbox}}
\newenvironment{gbadvertencia}
  {\begin{tcolorbox}[gbadmon,title={Advertencia}]\small}{\end{tcolorbox}}
\newenvironment{gbimportante}
  {\begin{tcolorbox}[gbadmon,title={Importante}]\small}{\end{tcolorbox}}
\newenvironment{gbprecaucion}
  {\begin{tcolorbox}[gbadmon,title={Precaución}]\small}{\end{tcolorbox}}

% ============================================================
% 8. CÓDIGO
% ------------------------------------------------------------
% listings Y NO minted: minted EXIGE --shell-escape Y Pygments,
% Y ESO ROMPE LA AUTONOMÍA DEL PROYECTO EN OTRA MÁQUINA.
% ============================================================
\lstset{
  basicstyle=\ttfamily\footnotesize,
  breaklines=true,
  columns=fullflexible,
  frame=none,
  showstringspaces=false,
  extendedchars=true,
  inputencoding=utf8
}

% ============================================================
% 9. COLOFÓN
% ------------------------------------------------------------
% LAS FUENTES Y LA VERSIÓN SON VALORES, NO TEXTO ESCRITO. SI
% UNA EDITORIAL REEMPLAZA LA ESTÉTICA Y CAMBIA LAS FUENTES,
% EL COLOFÓN SIGUE DICIENDO LA VERDAD SIN QUE NADIE SE
% ACUERDE DE EDITARLO.
% LA FECHA ES LA DE PUBLICACIÓN Y NO LA DEL DÍA: UNA MARCA DE
% TIEMPO HARÍA QUE CADA COMPILACIÓN PRODUJERA UN DIFF AUNQUE
% NADA HUBIERA CAMBIADO.
% ============================================================
% SI LA ESTÉTICA NO LAS DECLARA, EL COLOFÓN LO DICE EN VEZ DE DAR UN
% "Undefined control sequence" QUE NO EXPLICA NADA.
\providecommand{\gbfuentetexto}{(sin declarar)}
\providecommand{\gbfuentetitulos}{(sin declarar)}
\providecommand{\gbversion}{(sin version)}

% EL RECUADRO SE DEFINE COMO ESTILO Y NO EN LÍNEA: ASÍ UNA EDITORIAL
% QUE QUIERA OTRO ASPECTO REDEFINE gbcolofoncaja SIN TOCAR EL TEXTO,
% QUE ES LO QUE EL CONTRATO GARANTIZA.
\tcbset{gbcolofoncaja/.style={
  colback=black!7,
  colframe=black!50,
  boxrule=0.7pt,
  arc=3mm,
  left=8pt, right=8pt, top=8pt, bottom=8pt,
  boxsep=0pt}}

\newcommand{\gbTextoColofon}{%
  \begin{tcolorbox}[gbcolofoncaja]
    La composición tipográfica de este libro se realizó con
    gbpublisher versión~\gbversion, bajo un modelo de fuente única que
    garantiza la trazabilidad entre el original y cada una de sus
    salidas (HTML, ePub y PDF). Las familias tipográficas utilizadas son
    \gbfuentetexto\ para el texto e \gbfuentetitulos\ para los
    títulos. El proyecto se encuentra disponible en
    \url{https://github.com/albertomoyano/gbpublisher}.
  \end{tcolorbox}}

   entornos y macros del XSLT
%   % ============================================================
% preambulo-estetica.tex
% ------------------------------------------------------------
% CAPA REEMPLAZABLE. ES LA QUE UNA EDITORIAL CAMBIA POR LA SUYA
% CUANDO QUIERE MANTENER SU IDENTIDAD VISUAL. NO DEBE DEFINIR
% NINGUNO DE LOS ENTORNOS DEL CONTRATO.
% ------------------------------------------------------------
% ESTE ARCHIVO ES EL ESQUELETO. LOS BLOQUES DE DISEÑO QUE YA
% EXISTEN EN EL preambulo.tex EN USO —titlesec, titletoc,
% estilos de página, diseño de notas al pie, newfloat,
% imakeidx, glossaries, froufrou, siunitx— SE PEGAN EN EL
% APARTADO 8. NO ESTÁN ACÁ PORQUE SON DECISIONES DE DISEÑO
% PROPIAS Y NO CORRESPONDE REESCRIBIRLAS.
% ============================================================
% LAS CUATRO PRESTACIONES
% ------------------------------------------------------------
% El derivado emite PRESTACIONES y no un «modo». Así, agregar
% un estado editorial nuevo no obliga a tocar esta capa, y cada
% bandera dice qué hace en vez de cómo se llama.
%
%   estado_libro     showframe  homeoarchy  enlaces  cruces
%   borrador             si         si         si      no
%   en_correccion        no         si         si      no
%   en_produccion        no         no         si      no
%   publicado            no         no         no      si
%
% LAS CUATRO TOCAN OPCIONES DE PAQUETE, QUE NO SE PUEDEN
% SUPERPONER DESPUES DE CARGADO EL PAQUETE. POR ESO SON
% BANDERAS DEL DERIVADO Y NO REEMPLAZOS DE TEXTO.
% LA EXCEPCION ES hyperref: SE CARGA UNA SOLA VEZ EN EL
% CONTRATO Y LOS COLORES VAN POR \hypersetup, QUE SI SE
% SUPERPONE.
% ============================================================
% CORRECCIONES APLICADAS RESPECTO DEL PREAMBULO EN USO
%   - FUERA \usepackage[utf8]{luainputenc}: LuaLaTeX CONSUME
%     UTF-8 NATIVAMENTE (SC-02).
%   - FUERA \usepackage{svg}: EXIGE Inkscape Y --shell-escape,
%     Y EN LIBROS SOLO SE USAN PNG, JPG Y PDF.
%   - \renewcommand{\footnoterule} ESTABA DEFINIDO DOS VECES
%     CON EL MISMO CUERPO. QUEDA UNO.
%   - babel LO CARGA EL CONTRATO, ANTES DE csquotes Y biblatex,
%     QUE ES EL ORDEN QUE PIDE LA DOCUMENTACION.
%   - luatexja SOLO SI EL LIBRO LO NECESITA.
% ============================================================

\usepackage{ifthen}
\usepackage{calc}

% ============================================================
% 1. GEOMETRIA
% ------------------------------------------------------------
% Los valores vienen del derivado.
%
% showframe ES OPCION DE PAQUETE Y NO SE PUEDE SUPERPONER: SE
% PASA CON \PassOptionsToPackage PARA NO DUPLICAR LA LISTA
% ENTERA DE OPCIONES EN DOS RAMAS DE UN CONDICIONAL.
%
% Hoy los cuatro margenes son iguales en los 45 formatos de la
% base, asi que inner/outer da lo mismo que centering; se usan
% igual para que el dia que se carguen asimetricos no haya que
% tocar nada. En un libro encuadernado el margen interior suele
% necesitar mas aire que el exterior.
% ============================================================
\ifgbshowframe
  \PassOptionsToPackage{showframe}{geometry}
\fi

\usepackage[
  paperwidth=\gbanchopagina,
  paperheight=\gbaltopagina,
  inner=\gbmargeninterior,
  outer=\gbmargenexterior,
  top=\gbmargensuperior,
  bottom=\gbmargeninferior,
  includehead,
  includefoot,
  headsep=14pt,
  footskip=24pt
]{geometry}

% ------------------------------------------------------------
% CRUCES DE CORTE
% ------------------------------------------------------------
% La hoja de montaje es mayor que el formato final: el derivado
% la calcula sumando el margen de cruces por lado. Las cruces
% solo necesitan lugar, y ese lugar es el mismo para cualquier
% formato.
%
% crop DECLARA LA HOJA DE MONTAJE Y geometry EL FORMATO FINAL:
% SON DOS PARES DE MEDIDAS INDEPENDIENTES, NO UNA DERIVADA DE
% LA OTRA. VERIFICADO: crop EXPANDE MACROS EN SUS OPCIONES
% keyval, ASI QUE NO HACE FALTA \crop[...] APARTE.
% ------------------------------------------------------------
\ifgbcruces
  \usepackage[width=\gbanchomontaje,height=\gbaltomontaje,cam,center]{crop}
\fi

% ============================================================
% 2. IDIOMA
% ------------------------------------------------------------
% babel lo carga el contrato, antes de csquotes y biblatex.
% Aca solo se cambian las cadenas, que si son diseno.
% ============================================================
\addto\captionsspanish{\renewcommand{\contentsname}{Sumario}}

% ============================================================
% 3. TIPOGRAFIA
% ------------------------------------------------------------
% Las fuentes son de la estetica, no de la base. Por eso cuando
% una editorial reemplaza esta capa, cambia las fuentes y el
% colofon las nombra correctamente sin que haya que tocar nada
% mas: el contrato solo las lee.
% ============================================================
\usepackage{fontspec}
\usepackage[final,babel=true]{microtype}

\def\gbfuentetexto{Libertinus Serif}
\def\gbfuentetitulos{IBM Plex Sans Condensed}
\def\gbfuentemono{IBM Plex Mono}

\setmainfont{\gbfuentetexto}[
  Numbers={OldStyle,Proportional},
  Ligatures={TeX,Common,Discretionary},Scale=1.18]
\setsansfont{\gbfuentetitulos}[
Scale=MatchLowercase,
Ligatures=TeX,
Extension=.otf,
UprightFont=*-Regular,
ItalicFont=*-Italic,
BoldFont=*-SemiBold,
BoldItalicFont=*-SemiBoldItalic]
\setmonofont{\gbfuentemono}[
Scale=0.91,
Extension=.otf,
UprightFont=*-Regular,
ItalicFont = IBMPlexMono-Italic.otf,
BoldFont = IBMPlexMono-Bold.otf,
BoldItalicFont = IBMPlexMono-BoldItalic.otf]

%\setsansfont{\gbfuentetitulos}[Scale=MatchLowercase,Ligatures=TeX]
%\setmonofont{\gbfuentemono}[Scale=0.91]

\usepackage{unicode-math}
\setmathfont{Libertinus Math}[Scale=MatchLowercase]

% SOPORTE CJK SOLO SI EL LIBRO LO NECESITA: luatexja PESA EN
% CADA COMPILACION Y CASI NINGUN LIBRO LO USA.
\ifgbcjk
  \usepackage{luatexja-fontspec}
  \setmainjfont{FandolSong}
\fi

\renewcommand{\normalsize}{\fontsize{10pt}{14pt}\selectfont}
\normalsize

% ============================================================
% 4. LISTAS, NOTAS Y FILETES
% ============================================================
% REEMPLAZA EL APARTADO 4 DE preambulo-estetica.tex
%
% QUÉ CAMBIÓ RESPECTO DEL ESQUELETO, Y POR QUÉ
% ------------------------------------------------------------
% 1. SE SACÓ hang DE footmisc
%
%    hang SANGRA EL TEXTO DE LA NOTA HACIA ADENTRO PARA HACERLE
%    LUGAR A LA MARCA. LO QUE SE BUSCA ES LO CONTRARIO: LA MARCA
%    COLGADA EN EL MARGEN Y EL TEXTO OCUPANDO EL ANCHO COMPLETO
%    DE LA CAJA. ESO LO PRODUCE \deffootnote DE scrextend.
%
%    hang Y \deffootnote GOBIERNAN LA MISMA SANGRÍA POR
%    MECANISMOS DISTINTOS, ASÍ QUE CARGAR LOS DOS DEJA EL
%    RESULTADO A MERCED DE CUÁL SE APLIQUE ÚLTIMO. QUEDA UNO.
%
%    bottom Y stable SÍ SE CONSERVAN: FIJAN LAS NOTAS AL PIE DE
%    LA CAJA Y LAS PERMITEN EN LOS TÍTULOS, Y NINGUNO TOCA LA
%    SANGRÍA.
%
% 2. SE SACARON LOS DOS \patchcmd SOBRE \@footnotetext
%
%    VERIFICADO: LOS DOS FALLAN. scrextend REDEFINE
%    \@footnotetext, ASÍ QUE CUANDO EL PARCHE BUSCA \footnotesize
%    O \@MM ESOS TOKENS YA NO ESTÁN EN LA DEFINICIÓN.
%
%    EL PRIMERO SE REEMPLAZA POR \setkomafont{footnote}, QUE ES
%    EL MECANISMO PROPIO DE KOMA Y NO DEPENDE DE PARCHEAR NADA.
%    VERIFICADO: CON \small LA MISMA NOTA OCUPA DOS PÁGINAS Y CON
%    \footnotesize UNA.
%
%    EL SEGUNDO SE REEMPLAZA POR \interfootnotelinepenalty, EL
%    PARÁMETRO DEL KERNEL QUE GOBIERNA ESE CORTE. ADVERTENCIA
%    HONESTA: NO PUDE DEMOSTRAR QUE CAMBIE NADA. UNA NOTA DE
%    CATORCE PÁRRAFOS SE PARTIÓ IGUAL CON EL PARÁMETRO Y SIN ÉL.
%    QUEDA PORQUE ES EL MECANISMO DOCUMENTADO Y NO HACE DAÑO,
%    PERO SI EL COMPORTAMIENTO CON NOTAS LARGAS NO ES EL DEL
%    PREÁMBULO LEGACY, ESTE ES EL LUGAR A REVISAR.
% ============================================================

\usepackage{enumitem}
\setlist{nosep,topsep=4pt}

% ------------------------------------------------------------
% NOTAS AL PIE
% ------------------------------------------------------------
\usepackage[bottom,stable]{footmisc}
\usepackage{scrextend}

% SEPARACIÓN ENTRE NOTAS
\setlength{\footnotesep}{10pt}

% FILETE SEPARADOR: 0,5 pt DE GROSOR Y 40% DEL ANCHO DE COLUMNA
\renewcommand{\footnoterule}{%
	\kern-3pt
	\hrule height 0.5pt width 0.4\columnwidth
	\kern6pt}

% NUMERACIÓN: ENTRE CORCHETES, PALO SECO, CUERPO scriptsize
\renewcommand*{\thefootnote}{\scriptsize\sffamily{[\arabic{footnote}]}}

% CUERPO DE LA NOTA
\setkomafont{footnote}{\small}

% ------------------------------------------------------------
% MARCA COLGADA EN EL MARGEN
% ------------------------------------------------------------
% \llap SACA LA MARCA FUERA DE LA CAJA HACIA LA IZQUIERDA, ASÍ EL
% TEXTO DE LA NOTA EMPIEZA EN EL MARGEN Y OCUPA EL ANCHO COMPLETO.
%
% LOS TRES ARGUMENTOS DE \deffootnote SON: SANGRÍA DE LA PRIMERA
% LÍNEA, SANGRÍA DE LAS SIGUIENTES, Y CÓMO SE COMPONE LA MARCA.
% ------------------------------------------------------------
\newcommand*\gbmarcanotaespacio{0em}

\deffootnote{\gbmarcanotaespacio}
{\parindent}
{%
	\makebox[\gbmarcanotaespacio][r]{\llap{\thefootnotemark\quad}}%
}

% ------------------------------------------------------------
% NOTAS LARGAS PARTIDAS ENTRE PÁGINAS
% ------------------------------------------------------------
% Penalización de partir una nota entre dos páginas. Por omisión
% LaTeX pone 100 en la clase book, y el valor alto que usa el
% kernel en otros contextos hace que una nota muy larga desplace
% texto en lugar de partirse.
%
% VER LA ADVERTENCIA DEL ENCABEZADO: NO PUDE MEDIR DIFERENCIA.
% ------------------------------------------------------------
\interfootnotelinepenalty=100

% ------------------------------------------------------------
% RESTO DEL APARTADO
% ------------------------------------------------------------
\usepackage[labelfont=bf,font=small,labelsep=period,format=plain]{caption}
\usepackage{needspace}

\setcounter{tocdepth}{1}
\setcounter{secnumdepth}{3}
\renewcommand{\thepart}{\arabic{part}}

\raggedbottom
\clubpenalty=10000
\widowpenalty=10000

% ============================================================
% 5. MARCAS TIPOGRAFICAS DE TRABAJO
% ------------------------------------------------------------
% impnattypo marca homeoarquias y problemas de principio y fin
% de linea. La opcion draft es opcion de paquete y no se puede
% superponer: por eso la carga va condicionada y no comentada
% a mano.
% ============================================================
\ifgbhomeoarchy
  \IfFileExists{impnattypo.sty}{%
    \usepackage[hyphenation,homeoarchy,homeoarchywordcolor=orange,
                homeoarchycharcolor=orange,draft]{impnattypo}}{%
    \PackageWarning{gbpublisher}{impnattypo no instalado:
      se pierden las marcas tipograficas de trabajo}}
  \overfullrule=5pt
\fi

% ============================================================
% 6. COLOR DE LOS ENLACES
% ------------------------------------------------------------
% Solo el color. hyperref lo carga preambulo-hyperref.tex, que
% va después de esta capa y lee \ifgbenlaces para decidir entre
% enlaces coloreados y hidelinks.
% ============================================================
\usepackage{xcolor}
\definecolor{gbenlace}{RGB}{160,0,120}

% ============================================================
% 7. CONSTANCIA EN EL LOG
% ------------------------------------------------------------
% Deja escrito en el .log con que prestaciones se compuso.
% Cuando un PDF sale distinto de lo esperado, esta linea dice
% por que sin tener que abrir el derivado.
% ============================================================
\typeout{[gbpublisher] prestaciones:
  showframe=\ifgbshowframe si\else no\fi,
  homeoarchy=\ifgbhomeoarchy si\else no\fi,
  enlaces=\ifgbenlaces si\else no\fi,
  cruces=\ifgbcruces si\else no\fi}

% ============================================================
% 8. AQUI VAN LOS BLOQUES DE DISENO DEL PREAMBULO EN USO
% ------------------------------------------------------------
%   - titlesec y los estilos de titulo
%   - titletoc: partes, capitulos y secciones del sumario
%   - titleps: estilos de pagina y folios corridos
%     (\gbtituloabreviado ES LO QUE CORRESPONDE EN EL FOLIO
%      CORRIDO, Y HOY EL PREAMBULO NO LO CONSUME)
%   - newfloat: el entorno flotante "imagen"
%   - imakeidx + xindy: names (lo llena esindex desde biblatex),
%     onomastico y concepto (los carga el editor)
%   - glossaries: siglas y glosario
%   - siunitx, froufrou, qrcode, pdfpages, rotating
% ============================================================

% Cita con cambio de márgenes y tamaño tipografico
\renewenvironment{quote}
{\normalsize\list{}{\sf\leftmargin=14pt \rightmargin=0pt}%
	\item\relax}
{\endlist}

% Configuración de los índices
\usepackage[xindy]{imakeidx}
\makeindex[name=names, title={}, intoc=false]
\makeindex[name=onomastico, title={}, intoc=false]
\makeindex[name=concepto, title={}, intoc=false]


% ============================================================
% PERSONALIZACIÓN DE LOS ÍNDICES DE PARTES, CAPÍTULOS Y SECCIONES
% ------------------------------------------------------------
% Paquete requerido:
\usepackage{titletoc}
%
% OBJETIVO GENERAL:
% - Controlar la apariencia de los índices (Tabla de Contenidos)
%   para partes, capítulos y secciones.
% - Ajustar sangrías, tipografía, alineación y estilo de los
%   números de página.
% ============================================================

% ------------------------------------------------------------
% PARTES
% ------------------------------------------------------------
% \titlecontents{part}[<left>]{<above-code>}{<numbered-entry-format>}
%                        {<numberless-entry-format>}{<filler-page>}
%
% Argumentos usados:
%   [0em]                    -> Sangría izquierda de la entrada.
%   {\addvspace{5pt}\sf\bfseries\normalsize\selectfont\filright}
%                             -> Espacio vertical antes de la entrada (5pt),
%                                fuente sans serif, negrita, tamaño normal,
%                                alineación a la derecha (\filright).
%   {\contentslabel[\thecontentslabel]{2.5pc}}
%                             -> Formato de la etiqueta numérica (ej. Parte I)
%                                con ancho de 2.5pc.
%   {}                        -> Formato de entrada sin número.
%   {}                        -> No hay reglas ni página aquí.
\titlecontents{part}[0em]
{\addvspace{5pt}\sf\bfseries\normalsize\selectfont\filright}
{\contentslabel[\thecontentslabel]{2.5pc}}
{}
{}

% ------------------------------------------------------------
% CAPÍTULOS
% ------------------------------------------------------------
% Formato general:
%   - Sangría: 1.5pc
%   - Espacio vertical: 0.4em antes de cada capítulo
%   - Fuente: sans serif, alineación a la derecha (\filright)
%   - Número del capítulo alineado con margen de 1.5pc
%   - Separador de puntos hasta el número de página (\titlerule)
%
% Argumentos:
%   [1.5pc] -> sangría izquierda del capítulo
%   {\addvspace{.4em}\sf\selectfont\filright} -> formato general de la entrada
%   {\contentslabel{1.5pc}} -> ancho reservado para número
%   {\hspace*{-1.5pc}} -> corrección si no tiene número
%   {\titlerule*[1pc]{.}\contentspage[\hspace*{-4pc} {\rm\small\thecontentspage}]}
%      -> regla de puntos de relleno hasta el número de página,
%         página en tamaño pequeño y ligeramente desplazada
\titlecontents{chapter}[1.5pc]
{\addvspace{.4em}\sf\selectfont\filright}
{\contentslabel{1.5pc}}
{\hspace*{-1.5pc}}%
{\titlerule*[1pc]{.}\contentspage[\hspace*{-4pc} {\rm\small\thecontentspage}]}%
[]

% ------------------------------------------------------------
% SECCIONES
% ------------------------------------------------------------
% Formato general:
%   - Sangría: 4.5pc
%   - Fuente: \small, alineación a la derecha
%   - Número de sección: 2.5pc
%   - Ajuste de sangría si no hay número: -2.5pc
%   - Separador de puntos hasta el número de página (\titlerule)
\titlecontents{section}[4.5pc]
{\small\filright}                  % Formato general
{\contentslabel{2.5pc}}            % Número con ancho 2.5pc
{\hspace*{-2.5pc}}                 % Si no tiene número
{\titlerule*[1pc]{.}\contentspage} % Regla de puntos hasta la página


% ============================================================
% CONFIGURACIÓN DE ESTILOS DE TÍTULOS CON titlesec
% ------------------------------------------------------------
% Paquete requerido:
\usepackage[sf,bf,compact,topmarks,calcwidth,pagestyles,clearempty,newlinetospace]{titlesec}
%
% OBJETIVO GENERAL:
% - Personalizar la apariencia de partes, capítulos, secciones,
%   subsecciones, subsubsecciones, párrafos y subpárrafos.
% - Controlar fuente, tamaño, alineación, sangrías y espaciado.
% ============================================================

% ============================================================
% DISEÑO DE PARTES
% ------------------------------------------------------------
% Redefinición de \@part y \@spart para controlar:
% - Numeración de partes (si corresponde)
% - Entrada en el TOC
% - Centrado del título y subtítulo de la parte
% - Tamaño de fuente y estilo tipográfico
%
% \@part[#1]#2:
%   - #1: título corto para TOC y encabezado
%   - #2: título completo para mostrar en la página
%   - Centrado, fuente sans serif (\sf), \LARGE para nombre de parte
%     y \Large para título completo
%
% \@spart#1:
%   - Para partes no numeradas
%   - Solo muestra el título centrado, \LARGE, sans serif
%
\makeatletter
\def\@part[#1]#2{%
	\ifnum \c@secnumdepth >-2\relax
	\refstepcounter{part}%
	\addcontentsline{toc}{part}{Parte \thepart\hspace{1em}#1}%
	\else
	\addcontentsline{toc}{part}{#1}%
	\fi
	\markboth{}{}%
	{\centering
		\interlinepenalty \@M
		\ifnum \c@secnumdepth >-2\relax
		\sf\LARGE\selectfont \partname~\thepart
		\par\vskip 20\p@%
		\fi
		\sf\Large\selectfont
		#2\par}%
	\@endpart
}

\def\@spart#1{%
	\markboth{}{}%
	{\centering
		\interlinepenalty \@M
		\sf\LARGE\selectfont
		#1\par}%
	\@endpart
}
\makeatother

% ============================================================
% DISEÑO DE CAPÍTULO
% ------------------------------------------------------------
% \titleformat{\chapter}[display]{...}{...}{sep}{before-code}[after-code]
% - [display]: título en bloque separado (no inline)
% - \Large: tamaño base del título
% - \vspace{-2.5cm}: ajusta altura vertical (sube el título)
% - \centering: centrado
% - \textsc{\MakeLowercase\chaptertitlename}: 'Capítulo' en versalitas
% - \thechapter: número del capítulo
% - 1.5cm: separación entre número y texto
% - \filright\sf\LARGE: alineación a la derecha, fuente sans serif, tamaño LARGER
\titleformat{\chapter}[display]
{\Large}
{\centering{\textsc{\MakeLowercase\chaptertitlename}}~\thechapter}
{1.5cm}
{\filright\sf\LARGE}
[]
% ------------------------------------------------------------
% ESPACIADO DEL TÍTULO DE CAPÍTULO
% ------------------------------------------------------------
% \titlespacing*{\chapter}{izquierda}{antes}{después}
%
% ANTES = 0pt HACE QUE «capítulo N» ARRANQUE EN EL BORDE SUPERIOR
% DE LA CAJA. LA CLASE book RESERVA 50pt POR OMISIÓN.
%
% LA VERSIÓN CON ASTERISCO SUPRIME LA SANGRÍA DEL PRIMER PÁRRAFO
% DESPUÉS DEL TÍTULO.
%
% NO SE USA \vspace NEGATIVO DENTRO DEL \titleformat: ESO EMPUJA
% HACIA ARRIBA DESPUÉS DE QUE LaTeX YA RESERVÓ EL ESPACIO, Y EL
% RESULTADO DEPENDE DEL CONTEXTO. titlesec TIENE EL PARÁMETRO
% PARA DECIRLO DIRECTO.
% ------------------------------------------------------------
\titlespacing*{\chapter}{0pt}{0pt}{30pt}

% ============================================================
% DISEÑO DE SECCIÓN
% ------------------------------------------------------------
% Fuente: sans serif, negrita, tamaño grande (\large), alineación a la derecha
\titleformat{\section}[hang]
{\sf\bfseries\large\raggedright}
{\thesection}{.5em}{}[]
\titlespacing{\section}
{\parindent}{18pt plus 0.5pt minus 0.5pt}{6.75pt}

% ============================================================
% DISEÑO DE SUBSECCIÓN
% ------------------------------------------------------------
\titleformat{\subsection}[hang]
{\sf\large\raggedright}
{\thesubsection}{.5em}{}[]
\titlespacing{\subsection}
{\parindent}{18pt plus 0.5pt minus 0.5pt}{6.75pt}

% ============================================================
% DISEÑO DE SUBSUBSECCIÓN
% ------------------------------------------------------------
\titleformat{\subsubsection}[hang]
{\rm\bfseries\normalsize\raggedright}
{\thesubsubsection}{.5em}{}[]
\titlespacing{\subsubsection}
{\parindent}{18pt plus 0.5pt minus 0.5pt}{6.75pt}

% ============================================================
% DISEÑO DE PÁRRAFO
% ------------------------------------------------------------
% \paragraph: run-in (alineado inline, no bloque)
% - Negrita
% - Sin número
% - Raya al final del título (---)
\titlespacing{\paragraph}
{0pt}                % margen izquierda
{6.75pt plus 0.5pt minus 0.5pt} % espacio antes
{0pt}                % espacio después
\titleformat{\paragraph}[runin]
{\bfseries}{}{0em}{}[\mbox{ --- }]

% ============================================================
% DISEÑO DE SUBPÁRRAFO
% ------------------------------------------------------------
% \subparagraph: estilo colapsado, centrado
% - Fuente sans serif, negrita, tamaño \large
\titlespacing{\subparagraph}
{\parindent}{18pt plus 0.5pt minus 0.5pt}{6.75pt}
\titleformat{\subparagraph}[hang]
{\sf\bfseries\large\centering}
{\thesubparagraph}{.5em}{}[]

   diseño, reemplazable
%   % ============================================================
% preambulo-hyperref.tex
% ------------------------------------------------------------
% CUARTA Y ÚLTIMA CAPA DEL PREÁMBULO. VA DESPUÉS DE LA ESTÉTICA.
%
%   \input{preambulo-proyecto}   valores de la base
%   \input{preambulo-contrato}   entornos y macros del XSLT
%   \input{preambulo-estetica}   diseño, reemplazable
%   \input{preambulo-hyperref}   ESTA
%
% POR QUÉ ES UNA CAPA APARTE Y VA ÚLTIMA
% ------------------------------------------------------------
% hyperref PARCHEA MEDIO LaTeX Y TIENE QUE CARGARSE DESPUÉS DE
% TODO LO QUE TAMBIÉN PARCHEA: titlesec, titleps, footmisc,
% imakeidx. SI SE CARGA ANTES, ESOS PAQUETES LO PISAN Y LOS
% SÍNTOMAS SON LEJANOS Y DIFÍCILES DE ATRIBUIR.
%
% PERO biblatex NECESITA LO CONTRARIO: DETECTA hyperref AL
% CARGARSE, Y SI NO ESTÁ, LAS CITAS SALEN SIN ENLACE.
%
% LA SALIDA A ESA CONTRADICCIÓN NO ES MOVER NINGUNO DE LOS DOS:
% ES SEPARAR LA CARGA DE biblatex DEL REGISTRO DEL .bib.
%
%   preambulo-contrato   \input{cita-\estandar-config.tex}
%                        → carga biblatex, temprano
%   ESTA CAPA            \usepackage{hyperref}
%                        \addbibresource{...}
%                        → hyperref último, y el recurso
%                          bibliográfico DESPUÉS de hyperref
%
% ES EL ORDEN QUE YA ESTÁ VERIFICADO EN PRODUCCIÓN: ES EL DEL
% preambulo.tex DE LOS LIBROS COMPUESTOS EN LaTeX PURO, DONDE
% biblatex SE CARGA EN LA LÍNEA 669, hyperref EN LA 844 Y EL
% \addbibresource EN LA 858.
%
% TAMPOCO HACE FALTA \AtEndPreamble NI EXPANDIR MACROS A MANO:
% ESTA CAPA CORRE DESPUÉS DE LA ESTÉTICA, ASÍ QUE PUEDE LEER
% \ifgbenlaces Y \definecolor{gbenlace} DIRECTAMENTE.
% ============================================================

% ============================================================
% 1. OPCIONES SEGÚN LA PRESTACIÓN DE ENLACES
% ------------------------------------------------------------
% En la salida para imprenta los enlaces no deben imprimirse
% coloreados. hidelinks los deja del color del texto y sin
% recuadro, conservando la navegación dentro del PDF.
%
% SON OPCIONES DE PAQUETE Y NO \hypersetup POSTERIOR: hidelinks
% SOLO SURTE EFECTO COMO OPCIÓN DE CARGA.
% ============================================================
\ifgbenlaces
  \usepackage[
    unicode=true,
    pdfencoding=auto,
    psdextra,
    colorlinks=true,
    linkcolor=gbenlace,
    citecolor=gbenlace,
    urlcolor=gbenlace,
    filecolor=gbenlace,
    linktoc=all,
    breaklinks=true,
    bookmarks=true,
    bookmarksopen=true,
    bookmarksnumbered=true
  ]{hyperref}
\else
  \usepackage[
    unicode=true,
    pdfencoding=auto,
    psdextra,
    hidelinks,
    pdfborder={0 0 0},
    breaklinks=true,
    bookmarks=true,
    bookmarksopen=true,
    bookmarksnumbered=true
  ]{hyperref}
\fi

% ============================================================
% 2. METADATOS XMP
% ------------------------------------------------------------
% hyperref solo llena el diccionario Info heredado, que no tiene
% lugar para ISBN ni licencia. hyperxmp agrega el XMP con
% Dublin Core, que es lo que leen los repositorios.
%
% CADA CLAVE SE EMITE SOLO SI TIENE CONTENIDO: hyperxmp
% INSPECCIONA EL PRIMER CARÁCTER DE ALGUNOS VALORES Y, CON EL
% VALOR VACÍO, SE COME LA LLAVE DE CIERRE Y EL PARSER SE
% DESBANDA:
%   ! Argument of \hyxmp@first@char@i has an extra brace
%   ! File ended while scanning use of \KVS@@Process
%
% TÍTULO Y AUTORÍA VAN SIN GUARDA PORQUE EL DERIVADO LOS VALIDA
% ANTES DE ESCRIBIR: SI FALTAN, GenerarPreambuloProyecto ABORTA.
% ============================================================
\usepackage{hyperxmp}

\hypersetup{pdftitle={\gbtitulo}}
\hypersetup{pdfauthor={\gbautores}}
\ifdefempty{\gbsubtitulo}{}{\hypersetup{pdfsubject={\gbsubtitulo}}}
\ifdefempty{\gbpalabrasclave}{}{\hypersetup{pdfkeywords={\gbpalabrasclave}}}
\ifdefempty{\gbeditorial}{}{\hypersetup{pdfpublisher={\gbeditorial}}}
\ifdefempty{\gbanio}{}{\hypersetup{pdfdate={\gbanio}}}
\ifdefempty{\gblicencia}{}{\hypersetup{pdfcopyright={\gblicencia}}}
\ifdefempty{\gburllicencia}{}{\hypersetup{pdflicenseurl={\gburllicencia}}}

% ============================================================
% 3. RECURSO BIBLIOGRÁFICO
% ------------------------------------------------------------
% VA ACÁ Y NO JUNTO AL \usepackage DE biblatex: DESPUÉS DE
% hyperref. ES LA PIEZA QUE RESUELVE LA CONTRADICCIÓN DE ORDEN
% DESCRITA ARRIBA, Y ES ASÍ COMO ESTÁ EN EL PREÁMBULO EN
% PRODUCCIÓN DE LOS LIBROS EN LaTeX PURO.
%
% SI EL .bib NO EXISTE, \addbibresource ABORTA LA COMPILACIÓN.
% CON LA GUARDA, EL LIBRO SALE CON LAS CITAS SIN RESOLVER Y UN
% AVISO EN EL LOG, QUE ES MEJOR QUE NO SALIR: EL EDITOR VE EL
% PROBLEMA EN LAS PRUEBAS EN VEZ DE QUEDARSE SIN NADA QUE MIRAR.
% ============================================================
\IfFileExists{\gbarchivobib}{%
  \addbibresource{\gbarchivobib}}{%
  \PackageWarning{gbpublisher}{No se encontro el archivo de referencias: las citas saldran sin resolver}}
   ESTA
%
% POR QUÉ ES UNA CAPA APARTE Y VA ÚLTIMA
% ------------------------------------------------------------
% hyperref PARCHEA MEDIO LaTeX Y TIENE QUE CARGARSE DESPUÉS DE
% TODO LO QUE TAMBIÉN PARCHEA: titlesec, titleps, footmisc,
% imakeidx. SI SE CARGA ANTES, ESOS PAQUETES LO PISAN Y LOS
% SÍNTOMAS SON LEJANOS Y DIFÍCILES DE ATRIBUIR.
%
% PERO biblatex NECESITA LO CONTRARIO: DETECTA hyperref AL
% CARGARSE, Y SI NO ESTÁ, LAS CITAS SALEN SIN ENLACE.
%
% LA SALIDA A ESA CONTRADICCIÓN NO ES MOVER NINGUNO DE LOS DOS:
% ES SEPARAR LA CARGA DE biblatex DEL REGISTRO DEL .bib.
%
%   preambulo-contrato   \input{cita-\estandar-config.tex}
%                        → carga biblatex, temprano
%   ESTA CAPA            \usepackage{hyperref}
%                        \addbibresource{...}
%                        → hyperref último, y el recurso
%                          bibliográfico DESPUÉS de hyperref
%
% ES EL ORDEN QUE YA ESTÁ VERIFICADO EN PRODUCCIÓN: ES EL DEL
% preambulo.tex DE LOS LIBROS COMPUESTOS EN LaTeX PURO, DONDE
% biblatex SE CARGA EN LA LÍNEA 669, hyperref EN LA 844 Y EL
% \addbibresource EN LA 858.
%
% TAMPOCO HACE FALTA \AtEndPreamble NI EXPANDIR MACROS A MANO:
% ESTA CAPA CORRE DESPUÉS DE LA ESTÉTICA, ASÍ QUE PUEDE LEER
% \ifgbenlaces Y \definecolor{gbenlace} DIRECTAMENTE.
% ============================================================

% ============================================================
% 1. OPCIONES SEGÚN LA PRESTACIÓN DE ENLACES
% ------------------------------------------------------------
% En la salida para imprenta los enlaces no deben imprimirse
% coloreados. hidelinks los deja del color del texto y sin
% recuadro, conservando la navegación dentro del PDF.
%
% SON OPCIONES DE PAQUETE Y NO \hypersetup POSTERIOR: hidelinks
% SOLO SURTE EFECTO COMO OPCIÓN DE CARGA.
% ============================================================
\ifgbenlaces
  \usepackage[
    unicode=true,
    pdfencoding=auto,
    psdextra,
    colorlinks=true,
    linkcolor=gbenlace,
    citecolor=gbenlace,
    urlcolor=gbenlace,
    filecolor=gbenlace,
    linktoc=all,
    breaklinks=true,
    bookmarks=true,
    bookmarksopen=true,
    bookmarksnumbered=true
  ]{hyperref}
\else
  \usepackage[
    unicode=true,
    pdfencoding=auto,
    psdextra,
    hidelinks,
    pdfborder={0 0 0},
    breaklinks=true,
    bookmarks=true,
    bookmarksopen=true,
    bookmarksnumbered=true
  ]{hyperref}
\fi

% ============================================================
% 2. METADATOS XMP
% ------------------------------------------------------------
% hyperref solo llena el diccionario Info heredado, que no tiene
% lugar para ISBN ni licencia. hyperxmp agrega el XMP con
% Dublin Core, que es lo que leen los repositorios.
%
% CADA CLAVE SE EMITE SOLO SI TIENE CONTENIDO: hyperxmp
% INSPECCIONA EL PRIMER CARÁCTER DE ALGUNOS VALORES Y, CON EL
% VALOR VACÍO, SE COME LA LLAVE DE CIERRE Y EL PARSER SE
% DESBANDA:
%   ! Argument of \hyxmp@first@char@i has an extra brace
%   ! File ended while scanning use of \KVS@@Process
%
% TÍTULO Y AUTORÍA VAN SIN GUARDA PORQUE EL DERIVADO LOS VALIDA
% ANTES DE ESCRIBIR: SI FALTAN, GenerarPreambuloProyecto ABORTA.
% ============================================================
\usepackage{hyperxmp}

\hypersetup{pdftitle={\gbtitulo}}
\hypersetup{pdfauthor={\gbautores}}
\ifdefempty{\gbsubtitulo}{}{\hypersetup{pdfsubject={\gbsubtitulo}}}
\ifdefempty{\gbpalabrasclave}{}{\hypersetup{pdfkeywords={\gbpalabrasclave}}}
\ifdefempty{\gbeditorial}{}{\hypersetup{pdfpublisher={\gbeditorial}}}
\ifdefempty{\gbanio}{}{\hypersetup{pdfdate={\gbanio}}}
\ifdefempty{\gblicencia}{}{\hypersetup{pdfcopyright={\gblicencia}}}
\ifdefempty{\gburllicencia}{}{\hypersetup{pdflicenseurl={\gburllicencia}}}

% ============================================================
% 3. RECURSO BIBLIOGRÁFICO
% ------------------------------------------------------------
% VA ACÁ Y NO JUNTO AL \usepackage DE biblatex: DESPUÉS DE
% hyperref. ES LA PIEZA QUE RESUELVE LA CONTRADICCIÓN DE ORDEN
% DESCRITA ARRIBA, Y ES ASÍ COMO ESTÁ EN EL PREÁMBULO EN
% PRODUCCIÓN DE LOS LIBROS EN LaTeX PURO.
%
% SI EL .bib NO EXISTE, \addbibresource ABORTA LA COMPILACIÓN.
% CON LA GUARDA, EL LIBRO SALE CON LAS CITAS SIN RESOLVER Y UN
% AVISO EN EL LOG, QUE ES MEJOR QUE NO SALIR: EL EDITOR VE EL
% PROBLEMA EN LAS PRUEBAS EN VEZ DE QUEDARSE SIN NADA QUE MIRAR.
% ============================================================
\IfFileExists{\gbarchivobib}{%
  \addbibresource{\gbarchivobib}}{%
  \PackageWarning{gbpublisher}{No se encontro el archivo de referencias: las citas saldran sin resolver}}
   ESTA
%
% POR QUÉ ES UNA CAPA APARTE Y VA ÚLTIMA
% ------------------------------------------------------------
% hyperref PARCHEA MEDIO LaTeX Y TIENE QUE CARGARSE DESPUÉS DE
% TODO LO QUE TAMBIÉN PARCHEA: titlesec, titleps, footmisc,
% imakeidx. SI SE CARGA ANTES, ESOS PAQUETES LO PISAN Y LOS
% SÍNTOMAS SON LEJANOS Y DIFÍCILES DE ATRIBUIR.
%
% PERO biblatex NECESITA LO CONTRARIO: DETECTA hyperref AL
% CARGARSE, Y SI NO ESTÁ, LAS CITAS SALEN SIN ENLACE.
%
% LA SALIDA A ESA CONTRADICCIÓN NO ES MOVER NINGUNO DE LOS DOS:
% ES SEPARAR LA CARGA DE biblatex DEL REGISTRO DEL .bib.
%
%   preambulo-contrato   \input{cita-\estandar-config.tex}
%                        → carga biblatex, temprano
%   ESTA CAPA            \usepackage{hyperref}
%                        \addbibresource{...}
%                        → hyperref último, y el recurso
%                          bibliográfico DESPUÉS de hyperref
%
% ES EL ORDEN QUE YA ESTÁ VERIFICADO EN PRODUCCIÓN: ES EL DEL
% preambulo.tex DE LOS LIBROS COMPUESTOS EN LaTeX PURO, DONDE
% biblatex SE CARGA EN LA LÍNEA 669, hyperref EN LA 844 Y EL
% \addbibresource EN LA 858.
%
% TAMPOCO HACE FALTA \AtEndPreamble NI EXPANDIR MACROS A MANO:
% ESTA CAPA CORRE DESPUÉS DE LA ESTÉTICA, ASÍ QUE PUEDE LEER
% \ifgbenlaces Y \definecolor{gbenlace} DIRECTAMENTE.
% ============================================================

% ============================================================
% 1. OPCIONES SEGÚN LA PRESTACIÓN DE ENLACES
% ------------------------------------------------------------
% En la salida para imprenta los enlaces no deben imprimirse
% coloreados. hidelinks los deja del color del texto y sin
% recuadro, conservando la navegación dentro del PDF.
%
% SON OPCIONES DE PAQUETE Y NO \hypersetup POSTERIOR: hidelinks
% SOLO SURTE EFECTO COMO OPCIÓN DE CARGA.
% ============================================================
\ifgbenlaces
  \usepackage[
    unicode=true,
    pdfencoding=auto,
    psdextra,
    colorlinks=true,
    linkcolor=gbenlace,
    citecolor=gbenlace,
    urlcolor=gbenlace,
    filecolor=gbenlace,
    linktoc=all,
    breaklinks=true,
    bookmarks=true,
    bookmarksopen=true,
    bookmarksnumbered=true
  ]{hyperref}
\else
  \usepackage[
    unicode=true,
    pdfencoding=auto,
    psdextra,
    hidelinks,
    pdfborder={0 0 0},
    breaklinks=true,
    bookmarks=true,
    bookmarksopen=true,
    bookmarksnumbered=true
  ]{hyperref}
\fi

% ============================================================
% 2. METADATOS XMP
% ------------------------------------------------------------
% hyperref solo llena el diccionario Info heredado, que no tiene
% lugar para ISBN ni licencia. hyperxmp agrega el XMP con
% Dublin Core, que es lo que leen los repositorios.
%
% CADA CLAVE SE EMITE SOLO SI TIENE CONTENIDO: hyperxmp
% INSPECCIONA EL PRIMER CARÁCTER DE ALGUNOS VALORES Y, CON EL
% VALOR VACÍO, SE COME LA LLAVE DE CIERRE Y EL PARSER SE
% DESBANDA:
%   ! Argument of \hyxmp@first@char@i has an extra brace
%   ! File ended while scanning use of \KVS@@Process
%
% TÍTULO Y AUTORÍA VAN SIN GUARDA PORQUE EL DERIVADO LOS VALIDA
% ANTES DE ESCRIBIR: SI FALTAN, GenerarPreambuloProyecto ABORTA.
% ============================================================
\usepackage{hyperxmp}

\hypersetup{pdftitle={\gbtitulo}}
\hypersetup{pdfauthor={\gbautores}}
\ifdefempty{\gbsubtitulo}{}{\hypersetup{pdfsubject={\gbsubtitulo}}}
\ifdefempty{\gbpalabrasclave}{}{\hypersetup{pdfkeywords={\gbpalabrasclave}}}
\ifdefempty{\gbeditorial}{}{\hypersetup{pdfpublisher={\gbeditorial}}}
\ifdefempty{\gbanio}{}{\hypersetup{pdfdate={\gbanio}}}
\ifdefempty{\gblicencia}{}{\hypersetup{pdfcopyright={\gblicencia}}}
\ifdefempty{\gburllicencia}{}{\hypersetup{pdflicenseurl={\gburllicencia}}}

% ============================================================
% 3. RECURSO BIBLIOGRÁFICO
% ------------------------------------------------------------
% VA ACÁ Y NO JUNTO AL \usepackage DE biblatex: DESPUÉS DE
% hyperref. ES LA PIEZA QUE RESUELVE LA CONTRADICCIÓN DE ORDEN
% DESCRITA ARRIBA, Y ES ASÍ COMO ESTÁ EN EL PREÁMBULO EN
% PRODUCCIÓN DE LOS LIBROS EN LaTeX PURO.
%
% SI EL .bib NO EXISTE, \addbibresource ABORTA LA COMPILACIÓN.
% CON LA GUARDA, EL LIBRO SALE CON LAS CITAS SIN RESOLVER Y UN
% AVISO EN EL LOG, QUE ES MEJOR QUE NO SALIR: EL EDITOR VE EL
% PROBLEMA EN LAS PRUEBAS EN VEZ DE QUEDARSE SIN NADA QUE MIRAR.
% ============================================================
\IfFileExists{\gbarchivobib}{%
  \addbibresource{\gbarchivobib}}{%
  \PackageWarning{gbpublisher}{No se encontro el archivo de referencias: las citas saldran sin resolver}}
   ESTA
%
% POR QUÉ ES UNA CAPA APARTE Y VA ÚLTIMA
% ------------------------------------------------------------
% hyperref PARCHEA MEDIO LaTeX Y TIENE QUE CARGARSE DESPUÉS DE
% TODO LO QUE TAMBIÉN PARCHEA: titlesec, titleps, footmisc,
% imakeidx. SI SE CARGA ANTES, ESOS PAQUETES LO PISAN Y LOS
% SÍNTOMAS SON LEJANOS Y DIFÍCILES DE ATRIBUIR.
%
% PERO biblatex NECESITA LO CONTRARIO: DETECTA hyperref AL
% CARGARSE, Y SI NO ESTÁ, LAS CITAS SALEN SIN ENLACE.
%
% LA SALIDA A ESA CONTRADICCIÓN NO ES MOVER NINGUNO DE LOS DOS:
% ES SEPARAR LA CARGA DE biblatex DEL REGISTRO DEL .bib.
%
%   preambulo-contrato   \input{cita-\estandar-config.tex}
%                        → carga biblatex, temprano
%   ESTA CAPA            \usepackage{hyperref}
%                        \addbibresource{...}
%                        → hyperref último, y el recurso
%                          bibliográfico DESPUÉS de hyperref
%
% ES EL ORDEN QUE YA ESTÁ VERIFICADO EN PRODUCCIÓN: ES EL DEL
% preambulo.tex DE LOS LIBROS COMPUESTOS EN LaTeX PURO, DONDE
% biblatex SE CARGA EN LA LÍNEA 669, hyperref EN LA 844 Y EL
% \addbibresource EN LA 858.
%
% TAMPOCO HACE FALTA \AtEndPreamble NI EXPANDIR MACROS A MANO:
% ESTA CAPA CORRE DESPUÉS DE LA ESTÉTICA, ASÍ QUE PUEDE LEER
% \ifgbenlaces Y \definecolor{gbenlace} DIRECTAMENTE.
% ============================================================

% ============================================================
% 1. OPCIONES SEGÚN LA PRESTACIÓN DE ENLACES
% ------------------------------------------------------------
% En la salida para imprenta los enlaces no deben imprimirse
% coloreados. hidelinks los deja del color del texto y sin
% recuadro, conservando la navegación dentro del PDF.
%
% SON OPCIONES DE PAQUETE Y NO \hypersetup POSTERIOR: hidelinks
% SOLO SURTE EFECTO COMO OPCIÓN DE CARGA.
% ============================================================
\ifgbenlaces
  \usepackage[
    unicode=true,
    pdfencoding=auto,
    psdextra,
    colorlinks=true,
    linkcolor=gbenlace,
    citecolor=gbenlace,
    urlcolor=gbenlace,
    filecolor=gbenlace,
    linktoc=all,
    breaklinks=true,
    bookmarks=true,
    bookmarksopen=true,
    bookmarksnumbered=true
  ]{hyperref}
\else
  \usepackage[
    unicode=true,
    pdfencoding=auto,
    psdextra,
    hidelinks,
    pdfborder={0 0 0},
    breaklinks=true,
    bookmarks=true,
    bookmarksopen=true,
    bookmarksnumbered=true
  ]{hyperref}
\fi

% ============================================================
% 2. METADATOS XMP
% ------------------------------------------------------------
% hyperref solo llena el diccionario Info heredado, que no tiene
% lugar para ISBN ni licencia. hyperxmp agrega el XMP con
% Dublin Core, que es lo que leen los repositorios.
%
% CADA CLAVE SE EMITE SOLO SI TIENE CONTENIDO: hyperxmp
% INSPECCIONA EL PRIMER CARÁCTER DE ALGUNOS VALORES Y, CON EL
% VALOR VACÍO, SE COME LA LLAVE DE CIERRE Y EL PARSER SE
% DESBANDA:
%   ! Argument of \hyxmp@first@char@i has an extra brace
%   ! File ended while scanning use of \KVS@@Process
%
% TÍTULO Y AUTORÍA VAN SIN GUARDA PORQUE EL DERIVADO LOS VALIDA
% ANTES DE ESCRIBIR: SI FALTAN, GenerarPreambuloProyecto ABORTA.
% ============================================================
\usepackage{hyperxmp}

\hypersetup{pdftitle={\gbtitulo}}
\hypersetup{pdfauthor={\gbautores}}
\ifdefempty{\gbsubtitulo}{}{\hypersetup{pdfsubject={\gbsubtitulo}}}
\ifdefempty{\gbpalabrasclave}{}{\hypersetup{pdfkeywords={\gbpalabrasclave}}}
\ifdefempty{\gbeditorial}{}{\hypersetup{pdfpublisher={\gbeditorial}}}
\ifdefempty{\gbanio}{}{\hypersetup{pdfdate={\gbanio}}}
\ifdefempty{\gblicencia}{}{\hypersetup{pdfcopyright={\gblicencia}}}
\ifdefempty{\gburllicencia}{}{\hypersetup{pdflicenseurl={\gburllicencia}}}

% ============================================================
% 3. RECURSO BIBLIOGRÁFICO
% ------------------------------------------------------------
% VA ACÁ Y NO JUNTO AL \usepackage DE biblatex: DESPUÉS DE
% hyperref. ES LA PIEZA QUE RESUELVE LA CONTRADICCIÓN DE ORDEN
% DESCRITA ARRIBA, Y ES ASÍ COMO ESTÁ EN EL PREÁMBULO EN
% PRODUCCIÓN DE LOS LIBROS EN LaTeX PURO.
%
% SI EL .bib NO EXISTE, \addbibresource ABORTA LA COMPILACIÓN.
% CON LA GUARDA, EL LIBRO SALE CON LAS CITAS SIN RESOLVER Y UN
% AVISO EN EL LOG, QUE ES MEJOR QUE NO SALIR: EL EDITOR VE EL
% PROBLEMA EN LAS PRUEBAS EN VEZ DE QUEDARSE SIN NADA QUE MIRAR.
% ============================================================
\IfFileExists{\gbarchivobib}{%
  \addbibresource{\gbarchivobib}}{%
  \PackageWarning{gbpublisher}{No se encontro el archivo de referencias: las citas saldran sin resolver}}
